\chapter{The wave equation and Weyl asymptotics}\label{ch:wave}
\providecommand{\vol}{\operatorname{vol}}
\providecommand{\inj}{\operatorname{inj}}
\providecommand{\Lset}{\mathcal{L}}
\usetikzlibrary{patterns}

In this chapter the theory of conormal distributions (\cref{ch:conormal}) and
the Hamilton--Jacobi and transport equations (\cref{ch:hj}) are applied to
the wave operator $P=D_t^2-\Delta$ on $\R\times X$, where $X$ is a compact
Riemannian manifold. We first solve the Cauchy problem and construct the
forward fundamental solution, with energy estimates and finite speed of
propagation. Then we construct, for short times, a parametrix for the wave
kernel as a superposition of conormal distributions associated with
characteristic hypersurfaces (equivalently, as an oscillatory integral whose
phase solves the eikonal equation). Restricting this parametrix to the diagonal
and integrating gives the singularity at $t=0$ of the trace of the wave group,
$T(t)=\sum_j\cos(t\lambda_j)$; a Tauberian argument then yields Weyl's
asymptotic formula for the eigenvalues of $\Delta$ with remainder
$O(\lambda^{n-1})$. Finally we explain how the propagation of singularities
leads to the Poisson relation: the singularities of $T$ lie at $0$ and at the
lengths of closed geodesics. \Cref{wave:fig:roadmap} summarizes the logic.

\begin{figure}[ht]
  \centering
  \begin{tikzpicture}[
      box/.style={draw,rounded corners,align=center,font=\footnotesize,inner sep=3pt},>=Stealth]
    \node[box] (cauchy) at (0,0) {Cauchy problem by\\ eigenfunction expansion};
    \node[box] (speed) at (0,-1.6) {finite propagation\\ speed};
    \node[box] (param) at (4.6,-0.8) {short-time parametrix:\\ eikonal + transport};
    \node[box] (trace) at (4.6,-2.6) {singularity of $T(t)$\\ at $t=0$};
    \node[box] (weyl) at (9.6,-2.6) {Weyl's law\\ $N(\lambda)=C_W\lambda^n+O(\lambda^{n-1})$};
    \node[box] (prop) at (9.6,-0.8) {propagation of\\ singularities};
    \node[box] (poisson) at (9.6,-4.2) {Poisson relation};
    \draw[->] (cauchy) -- (param);
    \draw[->] (speed) -- (param);
    \draw[->] (param) -- (trace);
    \draw[->] (trace) -- node[below,font=\footnotesize]{Tauberian} (weyl);
    \draw[->] (param) -- (prop);
    \draw[->] (prop.east) to[bend left=40] (poisson.east);
    \draw[->] (trace) -- (poisson);
  \end{tikzpicture}
  \caption{Logic of the chapter. The analysis of the wave kernel for short times
  gives the leading singularity of the wave trace, and hence Weyl's law; propagation of
  singularities controls the wave trace for all times.}
  \label{wave:fig:roadmap}
\end{figure}

\begin{discussion}
Why should the wave equation say anything about eigenvalues? The wave trace
$T(t)=\sum_j\cos(t\lambda_j)$ is the Fourier transform of the \emph{evenized}
counting measure $\frac12\sum_j(\delta_{\lambda_j}+\delta_{-\lambda_j})$ (since
$\cos(t\lambda)=\frac12(e^{it\lambda}+e^{-it\lambda})$), which on $[0,\infty)$ carries the
same information as $\sum_j\delta_{\lambda_j}$. Large-$\lambda$ behaviour of the counting
measure corresponds to the singularities of $T$ in $t$, and the strongest
singularity sits at $t=0$. On the other hand, by finite propagation speed the wave
kernel for small $|t|$ only depends on the geometry of $X$ near the diagonal, where
it can be computed by a parametrix. Weyl's law comes from matching these two
descriptions of $T$ near $t=0$.
\end{discussion}

\section{Setup}\label{wave:sec:setup}

\subsection{The Laplacian on a compact manifold}

Throughout, $X$ is a compact $\CI$ manifold without boundary, of dimension
$n\ge1$, with a Riemannian metric $g$. In local coordinates
$g=\sum g_{ij}(x)\,dx^idx^j$, $(g^{ij})=(g_{ij})^{-1}$, $|g|=\det(g_{ij})$, the
Riemannian measure is $\dvol=|g|^{1/2}dx$ and the (non-negative) Laplacian is
\begin{equation*}
\Delta u=-|g|^{-1/2}\sum_{i,j=1}^n\pa_i\big(|g|^{1/2}g^{ij}\pa_ju\big),\qquad
\int_X\Delta u\,\bar v\dvol=\int_X\ang{du,dv}_g\dvol .
\end{equation*}
For a covector we write $|\xi|_{g(x)}=(\sum g^{ij}(x)\xi_i\xi_j)^{1/2}$ and set
$q(x,\xi)=|\xi|_{g(x)}$. We use $\dvol$ to identify densities with functions:
$\CmI(X)$ is the dual of $\CI(X)$ under $\pair{u}{v}=\int_Xuv\dvol$, and the
Schwartz kernel $K_A$ of an operator $A$ is taken with respect to $\dvol$,
$Au(x)=\int_XK_A(x,y)u(y)\dvol(y)$. In particular $K_{\Id}=|g(y)|^{-1/2}\delta(x-y)$
in local coordinates.

We use the following consequences of the elliptic theory on compact manifolds
(\cref{ch:manifolds}; see also \cite{Shubin}).
\begin{enumerate}[label=(F\arabic*)]
\item\label{wave:F1} There is an orthonormal basis $\{\phi_j\}_{j\ge0}$ of
  $L^2(X)$ of real eigenfunctions $\phi_j\in\CI(X)$, $\Delta\phi_j=\lambda_j^2\phi_j$,
  with $0=\lambda_0\le\lambda_1\le\lambda_2\le\dots\to\infty$ (repeated with
  multiplicity).
\item\label{wave:F2} For $s\in\R$, writing $u_j=\pair{u}{\phi_j}$ and
  $\ang{\lambda}=(1+\lambda^2)^{1/2}$,
  \begin{equation*}
  \Hs{s}(X)=\Big\{u\in\CmI(X);\ \norm{u}_s^2=\sum_j\ang{\lambda_j}^{2s}|u_j|^2<\infty\Big\}
  \end{equation*}
  with a norm equivalent to the usual one. Thus $\CI(X)=\bigcap_s\Hs{s}(X)$,
  $\CmI(X)=\bigcup_s\Hs{s}(X)$, every $u\in\CmI(X)$ is the sum of its expansion
  $\sum_ju_j\phi_j$ (in some $\Hs{s}$), and for $u\in\CI(X)$ the expansion
  converges in $\CI(X)$. For $s>n/2+k$ there is the Sobolev embedding
  $\norm{u}_{\mathcal C^k}\le C\norm{u}_s$.
\end{enumerate}
For a bounded Borel function $F$ on $[0,\infty)$ we set
$F(\Lambda)u=\sum_jF(\lambda_j)u_j\phi_j$; here $\Lambda=\Delta^{1/2}$. Such an
operator is bounded on every $\Hs{s}(X)$ with norm $\le\sup|F|$. We shall use
$\cos(t\Lambda)$ and $\Lambda^{-1}\sin(t\Lambda)$, the latter defined by the
function $\lambda\mapsto\sin(t\lambda)/\lambda$ (equal to $t$ at $\lambda=0$).
The counting function is
\begin{equation*}
N(\lambda)=\#\{j;\ \lambda_j\le\lambda\}
=\text{number of eigenvalues of $\Delta$ in $[0,\lambda^2]$, with multiplicity.}
\end{equation*}

\begin{lemma}\label{wave:lem:apriori}
For every $s>n/2$ there is $C_s$ with
$\sum_j\ang{\lambda_j}^{-2s}|\phi_j(x)|^2\le C_s^2$ for all $x\in X$; hence
$\sum_j\ang{\lambda_j}^{-2s}\le C_s^2\vol(X)$ and
$N(\lambda)\le C_s^2\vol(X)\ang{\lambda}^{2s}$.
\end{lemma}
\begin{proof}
By the Sobolev embedding, $u\mapsto u(x)$ is a continuous linear functional on
$\Hs{s}(X)$ of norm at most $C_s$, uniformly in $x$. In the spectral description
\ref{wave:F2}, $u(x)=\sum_ju_j\phi_j(x)$, and the norm of this functional is
exactly $(\sum_j\ang{\lambda_j}^{-2s}|\phi_j(x)|^2)^{1/2}$. Integrating over $X$
gives the second statement, and $N(\lambda)\ang{\lambda}^{-2s}\le\sum_{\lambda_j\le\lambda}\ang{\lambda_j}^{-2s}$
gives the third.
\end{proof}

In particular $\sum_j\ang{\lambda_j}^{-M}<\infty$ for $M>n$, and the sup norms
of $\phi_j$ and their derivatives grow at most polynomially in $\lambda_j$.

\subsection{The wave operator and its bicharacteristics}

The wave operator is $P=D_t^2-\Delta$ on $\R\times X$, $D_t=-i\pa_t$. In the
canonical coordinates $(t,x,\tau,\xi)$ of $T^*(\R\times X)=T^*\R\times T^*X$
its principal symbol is
\begin{equation*}
p(t,x,\tau,\xi)=\tau^2-|\xi|^2_{g(x)},
\end{equation*}
and the characteristic set is $\Sigma(P)=\Char(P)=\{p=0\}\setminus0$. On
$\Sigma(P)$ we have $\tau=\pm|\xi|_g$ with $\tau\ne0$, $\xi\ne0$, so
$d_{(\tau,\xi)}p=(2\tau,-2g^{-1}\xi)\ne0$: $P$ is of real principal type and
$\Sigma(P)$ is a smooth conic hypersurface. We use the Hamilton vector field
$\Ham{f}=\sum_j(\pa_{\xi_j}f\,\pa_{x_j}-\pa_{x_j}f\,\pa_{\xi_j})$ (with the
corresponding $\tau,t$ terms on $T^*(\R\times X)$); a \emph{null bicharacteristic}
of $P$ is an integral curve of $\Ham{p}$ in $\Sigma(P)$. Let
\begin{equation*}
G_s=\exp(s\Ham{q})\colon T^*X\setminus0\to T^*X\setminus0,\qquad q=|\xi|_g,
\end{equation*}
be the (homogeneous) geodesic flow. It is complete since $X$ is compact,
commutes with dilations in $\xi$ and preserves $q$, and the projection to $X$ of
$s\mapsto G_s(x,\xi)$ is the unit speed geodesic with initial velocity
$g^{-1}\xi/|\xi|_g$: indeed $\Ham{q}=q^{-1}\Ham{q^2/2}$ on $T^*X\setminus0$,
and Hamilton's equations for $\frac12q^2$ are the geodesic equations
($\dot x=g^{-1}\xi$ is the Legendre transform).

\begin{lemma}\label{wave:lem:bichar}
Along a null bicharacteristic, $\tau$ is constant and $t$ is strictly monotone.
Parametrised by $t$, the null bicharacteristic through $(\bar t,\bar x,\tau,\bar\xi)$ with
$\tau=\mp|\bar\xi|_g$ is
\begin{equation*}
t\longmapsto\big(t,G_{\pm(t-\bar t)}(\bar x,\bar\xi),\tau\big),\qquad t\in\R .
\end{equation*}
In particular its projection to $X$ is a unit speed geodesic.
\end{lemma}
\begin{proof}
$\Ham{p}=2\tau\pa_t-\Ham{q^2}=2\tau\pa_t-2q\Ham{q}$, so $\dot\tau=0$, $\dot t=2\tau\ne0$
and $d(x,\xi)/dt=-(q/\tau)\Ham{q}=\pm\Ham{q}$ when $\tau=\mp q$.
\end{proof}

\subsection{Conormal distributions and oscillatory integrals}

We recall from \cref{ch:conormal} the local description of conormal
distributions associated with a hypersurface, which is all we need. Let $M$ be
a manifold of dimension $N$ and $H\subset M$ a closed embedded hypersurface.
Then $u\in\Ic{m}(M,H)$ if $u\in\CmI(M)$ is smooth on $M\setminus H$ and near
each point of $H$, in local coordinates $(s,z)\in\R\times\R^{N-1}$ with
$H=\{s=0\}$,
\begin{equation}\label{wave:eq:conormal-local}
u=\frac1{2\pi}\int_\R e^{is\rho}a(z,\rho)\,d\rho+\CI,\qquad
a\in\Sym{m+\frac N4-\frac12}(\R^{N-1}_z;\R_\rho).
\end{equation}
(We work with functions rather than half-densities; in fixed coordinates the
two differ by a smooth non-vanishing factor.) The class of $a$ modulo symbols
of one order lower is the local form of the principal symbol. Any defining
function $h$ of $H$ (with $dh\ne0$ on $H$) may be used in place of $s$, since it
can be taken as a coordinate; and amplitudes $a(s,z,\rho)$ depending also on $s$
give elements of the same space (\cref{ch:conormal}; a special case is proved
in \cref{wave:lem:fourier-amplitude} below). For $M=\R$, $H=\{0\}$: a compactly
supported $u\in\CmI(\R)$ lies in $\Ic{m}(\R,\{0\})$ iff $\hat u\in\Sym{m-\frac14}(\R)$.
We always have $\WF(u)\subset N^*H$.

A hypersurface $H=\{h=0\}$ is \emph{characteristic} for $P$ if
$N^*H\subset\Sigma(P)$, i.e.\ $p(dh)=0$ on $H$; for $P$ this is the eikonal
equation $(\pa_th)^2=|d_xh|_g^2$. Characteristic hypersurfaces are obtained by
Hamilton--Jacobi theory, and the symbols of conormal solutions of $Pu\in\CI$
satisfy transport equations along $\Ham{p}$ (\cref{ch:hj}). For $P$ we carry
this out explicitly in \cref{wave:sec:parametrix}.

Finally we use standard facts about oscillatory integrals
\cite[Sections~8.1, 8.2]{Hormander-I}. Let $\Omega\subset\R^d$ be open and
$\Phi\in\CI(\Omega\times(\R^N\setminus0))$ real, homogeneous of degree one in
$\theta$, with $d_{z,\theta}\Phi\ne0$, and let $a\in\Sym{\mu}(\Omega;\R^N)$
vanish near $\theta=0$. Then
$I(z)=\int e^{i\Phi(z,\theta)}a(z,\theta)\,d\theta$, defined as
$\lim_{\eps\downarrow0}\int e^{i\Phi}a(z,\theta)\chi_0(\eps\theta)\,d\theta$ for any
$\chi_0\in\CIc(\R^N)$ with $\chi_0(0)=1$, is a distribution on $\Omega$, and
\begin{equation}\label{wave:eq:WF-osc}
\WF(I)\subset\{(z,d_z\Phi(z,\theta));\ d_\theta\Phi(z,\theta)=0,\ \theta\in\operatorname{cone\,supp}a\}.
\end{equation}
In particular $I\in\CI$ if $a\in\Sym{-\infty}$, or near points where
$d_\theta\Phi\ne0$ on the support of $a$. If some of the variables $z$ are
parameters $t$ and $d_{z',\theta}\Phi\ne0$ in the remaining variables, $I$ is a
smooth function of $t$ with values in distributions.

\section{The Cauchy problem}\label{wave:sec:cauchy}

\begin{theorem}\label{wave:thm:cauchy}
Let $s\in\R$, let $I\subset\R$ be an open interval containing $0$, and let
$u_0\in\Hs{s}(X)$, $u_1\in\Hs{s-1}(X)$, $f\in\Cz(I;\Hs{s-1}(X))$. There is a
unique $u\in\Cz(I;\Hs{s})\cap\mathcal C^1(I;\Hs{s-1})$ with
\begin{equation}\label{wave:eq:cauchy}
Pu=f\ \text{ in }I\times X,\qquad u(0)=u_0,\qquad D_tu(0)=u_1,
\end{equation}
namely
\begin{equation}\label{wave:eq:duhamel}
u(t)=\cos(t\Lambda)u_0+i\Lambda^{-1}\sin(t\Lambda)u_1-\int_0^t\Lambda^{-1}\sin((t-r)\Lambda)f(r)\,dr .
\end{equation}
It satisfies, for $t\in I$, the energy estimate
\begin{equation}\label{wave:eq:energy}
\norm{u(t)}_s+\norm{D_tu(t)}_{s-1}\le 4(1+|t|)\Big(\norm{u_0}_s+\norm{u_1}_{s-1}
+\Big|\int_0^t\norm{f(r)}_{s-1}dr\Big|\Big),
\end{equation}
and if $f=0$ the energy $\norm{\pa_tu(t)}^2_{L^2}+\norm{du(t)}^2_{L^2}$ is
independent of $t$ (for $s\ge1$). Uniqueness holds in the larger class
$\mathcal C^2(I;\Hs{-M}(X))$, any $M$. If $u_0,u_1\in\CI(X)$ and
$f\in\CI(I\times X)$ then $u\in\CI(I\times X)$.
\end{theorem}
\begin{proof}
Expanding in eigenfunctions, $u$ solves \eqref{wave:eq:cauchy} iff its
coefficients $u_j(t)=\pair{u(t)}{\phi_j}$ solve
$u_j''+\lambda_j^2u_j=-f_j$, $u_j(0)=u_{0,j}$, $u_j'(0)=iu_{1,j}$ (recall
$D_t^2=-\pa_t^2$ and $\pair{\Delta u}{\phi_j}=\lambda_j^2u_j$), whose unique
solution is the $j$-th coefficient of \eqref{wave:eq:duhamel}. This proves
uniqueness in $\mathcal C^2(I;\Hs{-M})$: if the data vanish, every $u_j$
vanishes, hence $u=0$ by \ref{wave:F2}. For the estimate use
$|\sin(t\lambda)/\lambda|\le\min(|t|,\lambda^{-1})\le2(1+|t|)\ang{\lambda}^{-1}$:
\begin{align*}
\ang{\lambda_j}^s|u_j(t)|&\le\ang{\lambda_j}^s|u_{0,j}|+2(1+|t|)\Big(\ang{\lambda_j}^{s-1}|u_{1,j}|
+\Big|\int_0^t\ang{\lambda_j}^{s-1}|f_j(r)|\,dr\Big|\Big),\\
\ang{\lambda_j}^{s-1}|D_tu_j(t)|&\le\ang{\lambda_j}^s|u_{0,j}|+\ang{\lambda_j}^{s-1}|u_{1,j}|
+\Big|\int_0^t\ang{\lambda_j}^{s-1}|f_j(r)|\,dr\Big|,
\end{align*}
the second because $D_tu_j=i\lambda_j\sin(t\lambda_j)u_{0,j}+\cos(t\lambda_j)u_{1,j}
+i\int_0^t\cos((t-r)\lambda_j)f_j(r)\,dr$. Taking $\ell^2$ norms in $j$ (and
Minkowski's inequality for the integral) gives \eqref{wave:eq:energy}; the same
bounds and dominated convergence show that \eqref{wave:eq:duhamel} lies in
$\Cz(I;\Hs{s})\cap\mathcal C^1(I;\Hs{s-1})$. If $f=0$ then
$|u_j'|^2+\lambda_j^2|u_j|^2$ is constant for each $j$, and summing gives the
conservation of $\norm{\pa_tu}^2+\pair{\Delta u}{\bar u}=\norm{\pa_tu}^2+\norm{du}^2$.
Smoothness is the case $Y=\{\mathrm{pt}\}$ of \cref{wave:prop:param}.
\end{proof}

The \emph{wave group} is the map
$U(t)\colon(u_0,u_1)\mapsto(u(t),D_tu(t))$ for $f=0$; explicitly
\begin{equation*}
U(t)=\begin{pmatrix}\cos t\Lambda&i\Lambda^{-1}\sin t\Lambda\\ i\Lambda\sin t\Lambda&\cos t\Lambda\end{pmatrix}.
\end{equation*}
By \cref{wave:thm:cauchy} it is a strongly continuous group of bounded operators
on $\Hs{s}(X)\oplus\Hs{s-1}(X)$ for each $s$: $U(t+t')=U(t)U(t')$ by
uniqueness. Its (distributional) trace is $2\sum_j\cos(t\lambda_j)$, see
\cref{wave:sec:trace}.

We need a version with parameters. Let $Y$ be a compact manifold (possibly a
point) and let $\Delta_X$ act on the first factor of $X\times Y$.

\begin{proposition}\label{wave:prop:param}
Let $I\ni0$ be an open interval. If $w\in\mathcal C^2(I;\Hs{-M}(X\times Y))$
satisfies $(D_t^2-\Delta_X)w=0$ and $w(0)=D_tw(0)=0$, then $w=0$. If
$F\in\CI(I\times X\times Y)$ and $w_0,w_1\in\CI(X\times Y)$, the solution of
$(D_t^2-\Delta_X)w=F$, $w(0)=w_0$, $D_tw(0)=w_1$, given by
\eqref{wave:eq:duhamel} with $\Lambda=\Lambda_X$, lies in $\CI(I\times X\times Y)$.
\end{proposition}
\begin{proof}
Choose a metric on $Y$ with Laplacian $\Delta_Y$ and eigenbasis $\psi_l$,
$\Delta_Y\psi_l=\mu_l^2\psi_l$. Then $\{\phi_j\otimes\psi_l\}$ is an orthonormal
basis of $L^2(X\times Y)$ of eigenfunctions of the Laplacian
$\Delta_X+\Delta_Y$ of the product metric, so by \ref{wave:F2} on $X\times Y$,
$\Hs{s}(X\times Y)$ is characterised by
$\sum_{j,l}(1+\lambda_j^2+\mu_l^2)^s|w_{jl}|^2<\infty$. Uniqueness follows as
before from the ODEs $w_{jl}''+\lambda_j^2w_{jl}=0$. The operators $\cos(t\Lambda_X)$
and $\Lambda_X^{-1}\sin(t\Lambda_X)$ act on the coefficients by multiplication by
functions bounded by $1+|t|$, so they are bounded on every $\Hs{s}(X\times Y)$,
locally uniformly and strongly continuously in $t$. Hence
$w\in\Cz(I;\Hs{s}(X\times Y))$ for all $s$; the equation gives
$\pa_t^2w=-\Delta_Xw-F\in\Cz(I;\Hs{s-2})$, and inductively
$w\in\mathcal C^k(I;\Hs{s}(X\times Y))$ for all $k,s$. By the Sobolev
embedding, $w\in\CI(I\times X\times Y)$.
\end{proof}

Before turning to singularities we show that waves travel at speed at most one.
The proof is an energy estimate on a shrinking ball (\cref{wave:fig:cone}): the energy
inside the ball $\{d(x_0,x)<R-t\}$ can only decrease with $t$, because energy can leave
the ball through its moving boundary but, travelling at unit speed, cannot enter.

\begin{figure}[ht]
  \centering
  \begin{tikzpicture}[>=Stealth,scale=1.0]
    \fill[gray!15] (-1,0) -- (-2.6,1.6) -- (2.6,1.6) -- (1,0) -- cycle;
    \fill[gray!15] (-1,0) -- (-2.6,-1.6) -- (2.6,-1.6) -- (1,0) -- cycle;
    \draw[dashed] (-1,0) -- (-2.6,1.6) (1,0) -- (2.6,1.6);
    \draw[dashed] (-1,0) -- (-2.6,-1.6) (1,0) -- (2.6,-1.6);
    \draw[->] (-4.2,0) -- (5.6,0) node[right]{$x$};
    \draw[->] (0,-1.8) -- (0,2.0) node[above]{$t$};
    \draw[line width=2.2pt] (-1,0) -- (1,0);
    \node[below] at (-0.5,0) {$K$};
    \node at (0,1.25) {$K_{|t|}$};
    \fill[pattern=north east lines] (2.6,0) -- (3.75,1.15) -- (4.9,0) -- cycle;
    \draw (2.6,0) -- (3.75,1.15) -- (4.9,0);
    \fill (3.75,0) circle (1.5pt) node[below]{$x_0$};
    \draw[<->] (3.75,-0.55) -- node[below,font=\small]{$R$} (4.9,-0.55);
    \node[right,font=\small,align=left] at (4.0,1.35) {$u=0$ on\\ $d(x_0,x)<R-t$};
  \end{tikzpicture}
  \caption{Finite propagation speed. If the Cauchy data vanish outside $K$, then $u(t)$
  vanishes outside $K_{|t|}$, the region between the dashed lines of slope $\pm1$. The
  proof shows that $u$ vanishes on each cone $\{d(x_0,x)<R-t\}$ (hatched) with base a
  ball disjoint from $K$.}
  \label{wave:fig:cone}
\end{figure}

\begin{theorem}[Finite propagation speed]\label{wave:thm:finite-speed}
Let $K\subset X$ be closed, $K_r=\{x\in X;\ d(x,K)\le r\}$ for $r\ge0$, where $d$ is
the Riemannian distance. If $u$ is the solution of \eqref{wave:eq:cauchy} with
$f=0$ and $\supp u_0\cup\supp u_1\subset K$ (any $s$), then
$\supp u(t)\subset K_{|t|}$ for all $t$.
\end{theorem}
\begin{proof}
First let $u_0,u_1\in\CI(X)$, so $u\in\CI(\R\times X)$. Fix
$x_0\in X$ and $0<R<\inj(X)$ (the injectivity radius, positive as $X$ is compact)
with $R<d(x_0,K)$. On the ball $B(x_0,\inj X)$ the function $d(x_0,\cdot)^2$ is
smooth, so $\rho_\eps=(d(x_0,\cdot)^2+\eps^2)^{1/2}$ is smooth there and
$|d\rho_\eps|_g=d(x_0,\cdot)/\rho_\eps<1$. Let $\theta_\delta\in\CI(\R)$ be
non-decreasing, $\theta_\delta=0$ on $(-\infty,0]$, $\theta_\delta=1$ on
$[\delta,\infty)$, and $\psi(t,x)=\theta_\delta(R-t-\rho_\eps(x))$, a smooth
function supported in $\{\rho_\eps\le R-t\}\subset B(x_0,R)$ for $t\ge0$. With
the energy density $e=|\pa_tu|^2+|du|_g^2$ and $\pa_t^2u=-\Delta u$ we have
$\pa_te=2\operatorname{Re}(\overline{\pa_tu}\,\pa_t^2u+\ang{d\pa_tu,du}_g)$ and, integrating by
parts, the terms with two derivatives of $u$ cancel:
\begin{equation*}
\frac{d}{dt}\int_X\psi e\dvol=\int_X(\pa_t\psi)e\dvol-2\operatorname{Re}\int_X\overline{\pa_tu}\,\ang{d\psi,du}_g\dvol
\le\int_X(\pa_t\psi+|d\psi|_g)e\dvol\le0,
\end{equation*}
because $\pa_t\psi=-\theta_\delta'$ and $|d\psi|_g=\theta_\delta'|d\rho_\eps|_g\le\theta_\delta'$.
At $t=0$ the integral vanishes since $u_0,u_1$ vanish on $\bar B(x_0,R)$.
Letting $\delta,\eps\to0$ we get $e(t,x)=0$ for $0\le t<R$,
$d(x_0,x)<R-t$, hence $\pa_tu=0$ there, and since $u(0)=0$ on $B(x_0,R)$,
$u=0$ on this cone. Given $0\le t<\inj(X)$ and $x_0\notin K_t$ we may choose
$R\in(t,\min(\inj X,d(x_0,K)))$, so $u(t,x_0)=0$: $\supp u(t)\subset K_t$ for
$0\le t<\inj X$. Since $(K_a)_b\subset K_{a+b}$ by the triangle inequality,
applying this repeatedly (to the solution with data at time $t_1$, using
uniqueness and time-translation invariance) gives all $t\ge0$; $t\le0$ follows
by time reversal $t\mapsto-t$.

For distributional data use duality: $\cos(t\Lambda)$ and $\Lambda^{-1}\sin(t\Lambda)$
are symmetric for the bilinear pairing (the $\phi_j$ are real). If
$\phi\in\CI(X)$ has $\supp\phi\cap K_{|t|}=\emptyset$, then
$d(\supp\phi,K)>|t|$, so by the smooth case $\cos(t\Lambda)\phi$ and
$\Lambda^{-1}\sin(t\Lambda)\phi$ vanish near $K$ and
$\pair{u(t)}{\phi}=\pair{u_0}{\cos(t\Lambda)\phi}+i\pair{u_1}{\Lambda^{-1}\sin(t\Lambda)\phi}=0$.
\end{proof}

\section{The forward fundamental solution}\label{wave:sec:forward}

For $\phi\in\CIc(\R\times X)$ set
\begin{equation*}
E\phi(t)=-\int_{-\infty}^t\Lambda^{-1}\sin((t-r)\Lambda)\phi(r)\,dr,\qquad
E^\dagger\phi(t)=-\int_t^\infty\Lambda^{-1}\sin((r-t)\Lambda)\phi(r)\,dr .
\end{equation*}
By \cref{wave:thm:cauchy} and \cref{wave:prop:param} (with data at a time
before, resp.\ after, $\supp\phi$) these are the unique solutions in $\CI(\R\times X)$ of
$Pu=\phi$ vanishing for $t\ll0$, resp.\ $t\gg0$, and they depend continuously on
$\phi$. The Schwartz kernel of $E$ is
\begin{equation}\label{wave:eq:E-kernel}
E(t,x;t',x')=-H(t-t')\,K_{\sin}(t-t',x,x'),\qquad
K_{\sin}(t,x,y)=\sum_j\frac{\sin(t\lambda_j)}{\lambda_j}\phi_j(x)\phi_j(y),
\end{equation}
($H$ the Heaviside function), and $E^\dagger$ is its transpose for the
pairing $\pair{u}{v}=\int uv\,dt\dvol$, since $K_{\sin}$ is symmetric in $(x,y)$.
For $A\subset\R\times X$ let
$J^\pm(A)=\{(t,x);\ \exists(t',x')\in A,\ d(x,x')\le\pm(t-t')\}$ be the causal
future and past.

\begin{theorem}\label{wave:thm:forward}
Let $f\in\CmI(\R\times X)$ have $\supp f\subset[T,\infty)\times X$ for some $T$.
Then there is a unique $u\in\CmI(\R\times X)$ with support bounded below in $t$
such that $Pu=f$. It satisfies $\supp u\subset J^+(\supp f)\subset[T,\infty)\times X$,
it is given by $\pair{u}{\phi}=\pair{f}{E^\dagger\phi}$ for
$\phi\in\CIc(\R\times X)$ (so $u=Ef$, extending $E$), and $u\in\CI$ if $f\in\CI$.
In particular $E$ is the unique forward fundamental solution of $P$:
$PE=EP=\Id$ on distributions with support bounded below, and
$\supp E\subset\{(t,x,t',x');\ d(x,x')\le t-t'\}$.
\end{theorem}
\begin{proof}
By \cref{wave:thm:finite-speed} (applied with initial time $\sup\{t;(t,x)\in\supp\phi\}$),
$\supp E^\dagger\phi\subset J^-(\supp\phi)$, so $\supp f\cap\supp E^\dagger\phi$ is
compact and $\pair{f}{E^\dagger\phi}$ is defined; it is continuous in $\phi$, so it
defines $u\in\CmI(\R\times X)$. Since $P$ is formally self-adjoint for the
bilinear pairing ($P^t=P$), $\pair{Pu}{\phi}=\pair{f}{E^\dagger P\phi}=\pair{f}{\phi}$,
because $E^\dagger P\phi-\phi$ solves the homogeneous equation and vanishes for $t\gg0$,
hence is zero by \cref{wave:thm:cauchy}. If $\supp\phi\cap J^+(\supp f)=\emptyset$
then $J^-(\supp\phi)\cap\supp f=\emptyset$, so $\pair{u}{\phi}=0$; this gives the
support statement. For uniqueness let $v$ have support bounded below and $Pv=0$.
For $\phi\in\CIc$, $\supp v\cap\supp E^\dagger\phi$ is compact, so
$\pair{v}{\phi}=\pair{v}{PE^\dagger\phi}=\pair{Pv}{E^\dagger\phi}=0$. If $f\in\CI$,
the Duhamel integral $-\int_{-\infty}^t\Lambda^{-1}\sin((t-r)\Lambda)f(r)\,dr$ is a smooth
solution (\cref{wave:prop:param}) vanishing for $t<T$, hence equals $u$. The
statements about $E$ are a reformulation, the support of the kernel coming from
\cref{wave:thm:finite-speed}.
\end{proof}

\begin{remark}\label{wave:rem:volterra}
On non-compact manifolds (for instance $\R^n$ with a metric that is Euclidean
outside a compact set) the spectral formula is not available, and one instead
constructs $E$ from a local parametrix $E'$ with $PE'=\Id+F$, where $F$ has a
smooth kernel supported in $t\ge t'$; $\Id+F$ is then inverted by the Volterra
(Neumann) series $\sum_k(-F)^{\circ k}$ for the $t$-convolution, which converges
because the $k$-fold composition gains a factor $|t-t'|^{k-1}/(k-1)!$ (see
\cref{wave:ex:volterra}). On a non-compact manifold the bound by $\vol(X)$ used in
that exercise must be replaced by uniform geometric control: one needs $F$ to be
uniformly bounded with kernel supported within a fixed distance of the diagonal
(finite propagation speed for the parametrix), and a uniform bound on the volumes of
balls of fixed radius; then each $z$-integration contributes at most
$\sup_x\vol(B(x,R))$ in place of $\vol(X)$, and the supports of the iterates grow
only linearly in $k$ over bounded times. This is only an outline and is not used
below.
\end{remark}

\section{The short-time structure of the wave kernel}\label{wave:sec:parametrix}

We now describe the Schwartz kernels
\begin{equation*}
K_{\cos}(t,x,y)=\sum_j\cos(t\lambda_j)\phi_j(x)\phi_j(y),\qquad K_{\sin}(t,x,y)=\sum_j\frac{\sin(t\lambda_j)}{\lambda_j}\phi_j(x)\phi_j(y)
\end{equation*}
of $\cos(t\Lambda)$ and $\Lambda^{-1}\sin(t\Lambda)$ for small $|t|$. For fixed
$y$, $K_{\cos}$ solves $PK_{\cos}=0$ with $K_{\cos}(0)=\delta_y$,
$D_tK_{\cos}(0)=0$, and $K_{\sin}$ solves $PK_{\sin}=0$ with $K_{\sin}(0)=0$,
$\pa_tK_{\sin}(0)=\delta_y$. For each $k\in\N_0$ both series converge in
$\mathcal C^k(\R;\Hs{-M_k}(X\times X))$ for $M_k$ large (depending on $k$), by
\cref{wave:lem:apriori}, since each $t$-derivative brings down a factor $\lambda_j$;
hence they converge in $\CI(\R;\CmI(X\times X))$. (No single $M$ works for all $k$:
on the circle the $2k$-th $t$-derivative of $K_{\cos}$ at $t=0$ has diagonal spectral
coefficients of size $|j|^{2k}$.)

\subsection{Plane waves}

To construct the wave kernel we must solve the Cauchy problem with the singular
datum $\delta_y$. The method of \cref{ch:hj} produces solutions conormal to
characteristic hypersurfaces, so we first write $\delta_y$ as a superposition of
distributions conormal to hyperplanes.

The initial datum $\delta_y$ is conormal to the point $\{y\}$, but the only
characteristic submanifolds for $P$ are hypersurfaces: if $Z\subset\R\times X$
has codimension $k$ and $N^*Z\subset\Sigma(P)\cup0$, then each fibre $N^*_zZ$ is a
$k$-dimensional linear space on which the Lorentzian quadratic form $p$ of
signature $(1,n)$ vanishes, so $k\le1$. The way around this is to write $\delta$
as a superposition of distributions conormal to hyperplanes. In $\R^n$, using
polar coordinates $\xi=r\omega$ and symmetrising $(r,\omega)\mapsto(-r,-\omega)$,
\begin{equation}\label{wave:eq:plane-waves}
\delta(x)=(2\pi)^{-n}\int_{\Sphere^{n-1}}\int_0^\infty e^{irx\cdot\omega}r^{n-1}\,dr\,d\omega
=\frac{1}{2(2\pi)^{n-1}}\int_{\Sphere^{n-1}}f_n(x\cdot\omega)\,d\omega,
\end{equation}
where
\begin{equation*}
f_n(s)=\frac1{2\pi}\int_\R e^{isr}|r|^{n-1}\,dr=|D_s|^{n-1}\delta(s).
\end{equation*}
For $n$ odd, $f_n=D_s^{n-1}\delta$; for $n$ even, $|D_s|^{n-1}$ is a
pseudodifferential operator and $f_n$ is not supported at $0$.

\begin{lemma}\label{wave:lem:fn}
$f_n\in\Ic{n-\frac34}(\R,\{0\})$ (modulo a smooth function), and the pull-back
$f_n(x\cdot\omega)$ under the submersion $\R^n\times\Sphere^{n-1}\ni(x,\omega)\mapsto x\cdot\omega$
lies in $\Ic{\frac n2-\frac14}(\R^n\times\Sphere^{n-1},R)$, $R=\{x\cdot\omega=0\}$.
\end{lemma}
\begin{proof}
Let $\rho\in\CI(\R)$ vanish near $0$ and equal $1$ for $|r|\ge1$. Then
$f_n=\frac1{2\pi}\int e^{isr}\rho(r)|r|^{n-1}dr+\frac1{2\pi}\int e^{isr}(1-\rho(r))|r|^{n-1}dr$;
the second term is smooth (the inverse Fourier transform of a compactly supported
function) and the first is of the form \eqref{wave:eq:conormal-local} with symbol
$\rho|r|^{n-1}\in\Sym{n-1}$, so of order $m$ with $m+\frac14-\frac12=n-1$. The map
$(x,\omega)\mapsto x\cdot\omega$ has differential $\omega\cdot dx+x\cdot d\omega\ne0$, so it
is a submersion, $R$ is a smooth hypersurface and $x\cdot\omega$ is a defining function.
The pull-back is $\frac1{2\pi}\int e^{i(x\cdot\omega)r}\rho(r)|r|^{n-1}dr+\CI$, with the same
symbol, now on a manifold of dimension $N=2n-1$; thus
$m+\frac{2n-1}4-\frac12=n-1$, i.e.\ $m=\frac n2-\frac14$.
\end{proof}

\begin{proposition}[Radon inversion]\label{wave:prop:radon}
If $\pi_1\colon\R^n\times\Sphere^{n-1}\to\R^n$ is the projection, then
$(\pi_1)_*\big(f_n(x\cdot\omega)\,d\omega\big)=2(2\pi)^{n-1}\delta$; this is \eqref{wave:eq:plane-waves}.
\end{proposition}
\begin{proof}
For $\phi\in\cS(\R^n)$ and fixed $\omega$, $\int f_n(x\cdot\omega)\phi(x)\,dx=\pair{f_n}{\mathcal R_\omega\phi}$,
where $\mathcal R_\omega\phi(s)=\int_{x\cdot\omega=s}\phi$ satisfies
$\widehat{\mathcal R_\omega\phi}(r)=\hat\phi(r\omega)$. Hence
$\pair{f_n}{\mathcal R_\omega\phi}=\frac1{2\pi}\int_\R|r|^{n-1}\hat\phi(-r\omega)\,dr$ and, by
polar coordinates and Fourier inversion,
$\int_{\Sphere^{n-1}}\pair{f_n}{\mathcal R_\omega\phi}\,d\omega=\frac2{2\pi}\int_{\R^n}\hat\phi(\xi)\,d\xi=2(2\pi)^{n-1}\phi(0)$.
\end{proof}

Thus, to solve the Cauchy problem with datum $\delta_y$ we may solve it, for each
$\omega$, with data conormal to the hyperplane $\{(x-y)\cdot\omega=0\}$, using the two
characteristic hypersurfaces through it, and then integrate over $\omega$. It is
convenient to keep the radial variable as well: this amounts to an oscillatory
integral whose phase solves the eikonal equation.

\begin{question}
Why are there two characteristic hypersurfaces through each hyperplane, and why do we
need both?
\end{question}
\begin{answer}
Over a conormal $\omega\,dx$ to the hyperplane the characteristic set contains the
two points $\tau=\pm|\omega|_g$, giving two families of bicharacteristics and hence two
characteristic hypersurfaces: a plane wave moving in the direction $\omega$ and one
moving in the direction $-\omega$. The wave equation is of second order in $t$, so we
must match two pieces of Cauchy data, $u(0)$ and $D_tu(0)$; a single travelling wave
has its time derivative determined by its value, and only the sum of the two
can match arbitrary data (\cref{wave:ex:initial-split}). In one dimension this is
d'Alembert's formula $\cos(tD_x)\delta=\frac12(\delta(x-t)+\delta(x+t))$.
\end{answer}

\subsection{The eikonal equation}

Work in a coordinate chart $U\subset\R^n$ of $X$, and let $q(x,\eta)=|\eta|_{g(x)}$.

\begin{lemma}\label{wave:lem:eikonal}
Let $W\Subset U$ be open. There exist $t_0>0$ and real
$\varphi_\pm\in\CI((-t_0,t_0)\times W\times(\R^n\setminus0))$, homogeneous of degree
one in $\eta$, with
\begin{equation}\label{wave:eq:eikonal}
\pa_t\varphi_\pm=\pm q(x,d_x\varphi_\pm),\qquad \varphi_\pm(0,x,\eta)=x\cdot\eta .
\end{equation}
They are unique in the following sense: if $W\Subset U'\subset U$, $U'$ open, and
$\tilde\varphi_\pm$ are real solutions of \eqref{wave:eq:eikonal} on
$(-t_1,t_1)\times U'\times(\R^n\setminus0)$, then $\tilde\varphi_\pm=\varphi_\pm$ on
$(-t_2,t_2)\times W\times(\R^n\setminus0)$, $t_2=\min(t_0,t_1,\operatorname{dist}(W,
X\setminus U'))$. (Initial data on $W$ alone do not suffice: for $n=1$, flat metric,
$W=(0,1)$ and $\eta>0$, both $\eta(x+t)$ and $\eta F(x+t)$, with $F(s)=s$ for $s\le1$,
$F'>0$ and $F(s)\ne s$ for $1<s<1+\delta$, solve the equation with
sign $+$ and agree at $t=0$ on $W$, but differ at points of $(0,t)\times W$ for every
$t>0$.)
Moreover $\varphi_-(t,x,\eta)=\varphi_+(-t,x,\eta)$, $d_x\varphi_\pm\ne0$, and for $y\in W$:
\begin{equation}\label{wave:eq:critical}
d_\eta\big(\varphi_\pm(t,x,\eta)-y\cdot\eta\big)=0\iff\big(x,d_x\varphi_\pm(t,x,\eta)\big)=G_{\mp t}(y,\eta),
\end{equation}
in which case $(x,d_x\varphi_\pm)$ lies on the geodesic flow-out and
$d(x,y)\le|t|$.
\end{lemma}
\begin{proof}
This is the method of characteristics (the Hamilton--Jacobi theorem of
\cref{ch:hj}) for $F_\pm=\tau\mp q(x,\xi)$ with initial surface $\{t=0\}$, which is
non-characteristic since $\pa_\tau F_\pm=1$. Explicitly, consider $\varphi_+$ and let
$(x(t),\xi(t))=G_{-t}(x_0,\eta)$, i.e.\ $\dot x=-\pa_\xi q$, $\dot\xi=\pa_xq$, which are the
$x,\xi$ components of the bicharacteristics of $F_+$. For $|t|<t_0$ small, uniformly for
$|\eta|=1$ and $x_0$ near $\bar W$, the map $x_0\mapsto x(t;x_0,\eta)$ is a
diffeomorphism onto a neighbourhood of $\bar W$, and $\varphi_+$ is defined by
$\varphi_+(t,x(t;x_0,\eta),\eta)=x_0\cdot\eta$, extended homogeneously; the Hamilton--Jacobi
theorem shows that it solves \eqref{wave:eq:eikonal} and that
$d_x\varphi_+(t,x(t;x_0,\eta),\eta)=\xi(t)\ne0$ (consistently, $\frac d{dt}\varphi_+(t,x(t),\eta)
=q-\xi\cdot\pa_\xi q=0$ by Euler's relation). Uniqueness also follows from the
characteristics: for a solution $\tilde\varphi_+$ on $(-t_1,t_1)\times U'$, the curves
$\dot x=-\pa_\xi q(x,d_x\tilde\varphi_+)$ have unit speed, so the one through $(t,x)$,
$x\in W$, $|t|<t_2$, stays in $U'$ back to time $0$; along it $\tilde\varphi_+$ is constant
and $(x,d_x\tilde\varphi_+)$ satisfies Hamilton's equations, so $(x,d_x\tilde\varphi_+)$ is
the bicharacteristic $G_{-t}(x_0,\eta)$ with the same initial point as for $\varphi_+$, and
$\tilde\varphi_+(t,x,\eta)=x_0\cdot\eta=\varphi_+(t,x,\eta)$. Since $\psi(t)=\varphi_+(-t)$ solves the equation with the sign $-$,
$\varphi_-(t)=\varphi_+(-t)$.

For \eqref{wave:eq:critical}, differentiate $\varphi_+(t,x(t;x_0,\eta),\eta)=x_0\cdot\eta$ in
$\eta$:
\begin{equation*}
d_\eta\varphi_+(t,x(t),\eta)+\xi(t)\cdot\pa_\eta x(t;x_0,\eta)=x_0 .
\end{equation*}
The second term vanishes: with $w=\pa_\eta x\cdot\delta\eta$, $v=\pa_\eta\xi\cdot\delta\eta$,
$\frac d{dt}(\xi\cdot w)=\pa_xq\cdot w-\xi\cdot(\pa_x\pa_\xi q\,w+\pa_\xi^2q\,v)=0$, since
$\xi\cdot\pa_\xi\pa_xq=\pa_xq$ and $\xi\cdot\pa_\xi^2q=0$ by homogeneity, and $w=0$ at $t=0$.
(This is the invariance of the Liouville form $\xi\cdot dx$ under the homogeneous
flow $G_s$.) Hence $d_\eta\varphi_+(t,x,\eta)=y$ iff $x=x(t;y,\eta)$, i.e.\ iff
$(x,d_x\varphi_+(t,x,\eta))=G_{-t}(y,\eta)$. The case of $\varphi_-$ follows from
$\varphi_-(t)=\varphi_+(-t)$. Since $G_s$ moves base points along unit speed
geodesics, $d(x,y)\le|t|$.
\end{proof}

\subsection{Transport equations}

For a real phase $\phi(t,x)$ and an amplitude $a(t,x)$, a direct computation
(using $\Delta e^{i\phi}=e^{i\phi}(|d_x\phi|_g^2+i\Delta\phi)$) gives
\begin{equation}\label{wave:eq:P-on-exp}
P(e^{i\phi}a)=e^{i\phi}\Big[\big((\pa_t\phi)^2-|d_x\phi|_g^2\big)a-iL_\phi a+Pa\Big],
\end{equation}
\begin{equation*}
L_\phi a=2\big(\pa_t\phi\,\pa_ta-\ang{d_x\phi,d_xa}_g\big)+\big(\pa_t^2\phi+\Delta\phi\big)a .
\end{equation*}
If $\phi$ solves the eikonal equation the first term vanishes, and $L_\phi$ is a
first order operator whose principal part is the vector field
$V_\phi=2(\pa_t\phi\,\pa_t-\sum g^{ij}\pa_i\phi\,\pa_j)$, the projection to $\R\times X$
of $\Ham{p}$ restricted to the Lagrangian manifold $\{(t,x,d\phi)\}$. For
$\phi=\varphi_\pm(t,x,\eta)-y\cdot\eta$ the $\pa_t$-coefficient is $\pm2q(x,d_x\varphi_\pm)\ne0$,
so $V_\phi$ is transversal to $\{t=0\}$ and its integral curves are (reparametrised)
unit speed geodesics. If $a\sim\sum_{k\ge0}a_k$ with $a_k$ homogeneous of degree
$\mu-k$ in $\eta$, then $P(e^{i\phi}a)=e^{i\phi}r$ with $r\in\Sym{-\infty}$ in $\eta$
provided the \emph{transport equations}
\begin{equation}\label{wave:eq:transport}
L_\phi a_0=0,\qquad L_\phi a_k=-iPa_{k-1}\quad(k\ge1)
\end{equation}
hold; both sides of the $k$-th equation are homogeneous of degree $\mu+1-k$. Each is
a linear ODE along the integral curves of $V_\phi$, uniquely solvable for small $|t|$
given $a_k(0,\cdot)$, and the solution is homogeneous of the right degree.

\subsection{The parametrix}

Fix $\chi\in\CI(\R^n)$ with $\chi(\eta)=0$ for $|\eta|\le\frac12$ and $\chi(\eta)=1$ for
$|\eta|\ge1$. Choose a finite cover of $X$ by coordinate charts $U_k$, open sets
$V_k\Subset U_k$ still covering $X$, a partition of unity $\chi_k$ with
$\supp\chi_k\subset V_k$, and $\delta_0>0$ such that the $3\delta_0$-neighbourhood of
$\bar V_k$ is relatively compact in $U_k$ for all $k$. Let $W_k$ be the
$2\delta_0$-neighbourhood of $\bar V_k$ and $\zeta_k\in\CIc(W_k)$ with $\zeta_k=1$ on the
$\delta_0$-neighbourhood of $\bar V_k$.

\begin{theorem}\label{wave:thm:parametrix}
There are $t_0\in(0,\delta_0)$, phase functions $\varphi_{k,\pm}$ on
$(-t_0,t_0)\times W_k\times(\R^n\setminus0)$ as in \cref{wave:lem:eikonal}, and
polyhomogeneous symbols $a_{k,\pm}\in\Symph{0}$, $b_{k,\pm}\in\Symph{-1}$ in
$\eta\in\R^n$, depending smoothly on $(t,x,y)\in(-t_0,t_0)\times W_k\times V_k$, such that
\begin{align*}
K_{\cos}(t,x,y)&\equiv\sum_{k,\pm}\zeta_k(x)\chi_k(y)(2\pi)^{-n}\int e^{i(\varphi_{k,\pm}(t,x,\eta)-y\cdot\eta)}a_{k,\pm}(t,x,y,\eta)\chi(\eta)\,d\eta,\\
K_{\sin}(t,x,y)&\equiv\sum_{k,\pm}\zeta_k(x)\chi_k(y)(2\pi)^{-n}\int e^{i(\varphi_{k,\pm}(t,x,\eta)-y\cdot\eta)}b_{k,\pm}(t,x,y,\eta)\chi(\eta)\,d\eta,
\end{align*}
modulo $\CI((-t_0,t_0)\times X\times X)$. At $t=0$,
\begin{equation}\label{wave:eq:initial-symbols}
a_{k,\pm}(0,x,y,\eta)-\tfrac12|g(y)|^{-1/2}\in\Sym{-1},\qquad
b_{k,\pm}(0,x,y,\eta)\mp\frac{|g(y)|^{-1/2}}{2i\,q(x,\eta)}\in\Sym{-2}.
\end{equation}
\end{theorem}
\begin{proof}
Apply \cref{wave:lem:eikonal} on $W_k$ (with $W_k\Subset U_k$), shrinking $t_0$
so that also all integral curves of the transport fields starting in $W_k$ stay in
the $3\delta_0$-neighbourhood for $|t|<t_0$. We construct the cosine term; write
$\phi_\pm=\varphi_{k,\pm}-y\cdot\eta$ and $S=\sum_\pm(2\pi)^{-n}\int e^{i\phi_\pm}a_\pm\chi\,d\eta$.

\emph{Amplitudes.} We seek $a_\pm\sim\sum_{j\ge0}a_{\pm,j}$, $a_{\pm,j}$ homogeneous
of degree $-j$, solving \eqref{wave:eq:transport} for $\phi_\pm$, with initial values
fixed by the requirements $S(0)\equiv|g(y)|^{-1/2}\delta(x-y)$, $D_tS(0)\equiv0$.
Since $\phi_\pm(0)=(x-y)\cdot\eta$ and $\pa_t\phi_\pm(0)=\pm q(x,\eta)$, these amount to
\begin{gather*}
a_{+,0}+a_{-,0}=|g(y)|^{-1/2},\qquad a_{+,j}+a_{-,j}=0\ (j\ge1),\\
q(x,\eta)(a_{+,j}-a_{-,j})+D_ta_{+,j-1}+D_ta_{-,j-1}=0\ (j\ge0)
\end{gather*}
at $t=0$ (with $a_{\pm,-1}=0$), the last being the part homogeneous of degree
$1-j$ of $\sum_\pm(\pa_t\phi_\pm a_\pm+D_ta_\pm)|_{t=0}$. Inductively in $j$: the
initial values of $a_{\pm,j}$ are determined by these equations, given
$a_{\pm,j-1}$ for all $t$, and then $a_{\pm,j}$ is determined for $|t|<t_0$ by the
transport equation. In particular $a_{\pm,0}(0)=\frac12|g(y)|^{-1/2}$. Let $a_\pm$ be
a Borel sum of the series, multiplied by $\chi$. Then
$PS=(2\pi)^{-n}\sum_\pm\int e^{i\phi_\pm}r_\pm\,d\eta$ with $r_\pm\in\Sym{-\infty}$, so
$PS\in\CI$; similarly $S(0)-|g(y)|^{-1/2}\delta(x-y)$ and $D_tS(0)$ are oscillatory
integrals with amplitudes in $\Sym{-\infty}$, hence smooth. For the sine term the
conditions are $\sum_\pm b_\pm(0)\in\Sym{-\infty}$ and
$\sum_\pm(i\pa_t\phi_\pm b_\pm+\pa_tb_\pm)(0)-|g(y)|^{-1/2}\in\Sym{-\infty}$, giving
$b_{\pm,0}(0)=\pm|g(y)|^{-1/2}/(2iq(x,\eta))$ and then the same induction.

\emph{Singular support.} The phase $\phi_\pm$ is non-degenerate in $(y,\eta)$
($d_y\phi_\pm=-\eta\ne0$), so $S$ is a smooth function of $t$ with values in
distributions on $W_k\times V_k$, and by \eqref{wave:eq:WF-osc} and
\eqref{wave:eq:critical} it is smooth on $\{d(x,y)>|t|\}$.

\emph{Comparison.} Let $w=\chi_k(y)K_{\cos}-\zeta_k(x)\chi_k(y)S$ on
$(-t_0,t_0)\times X\times X$. Then
\begin{equation*}
Pw=-\zeta_k\chi_kPS-[P,\zeta_k]\chi_kS\in\CI,
\end{equation*}
the commutator term because $d\zeta_k=0$ on the $\delta_0$-neighbourhood of
$\supp\chi_k$, while $S$ is smooth where $d(x,y)>|t|$ and $|t|<t_0<\delta_0$. Also
$w(0)=\chi_k(y)(K_{\Id}-\zeta_k(x)S(0))$ and $D_tw(0)$ are smooth, since
$\zeta_k(x)=1$ near the support of $\chi_k(y)\delta(x-y)$. Both $K_{\cos}$ and $S$ lie
in $\mathcal C^2((-t_0,t_0);\Hs{-M}(X\times X))$ for $M$ large (only two time
derivatives are needed, and each costs a fixed number of Sobolev orders), so by
\cref{wave:prop:param}
(with $Y=X$ the $y$-factor) $w\in\CI$. Summing over $k$ proves the claim for
$K_{\cos}$, and the same argument applies to $K_{\sin}$.
\end{proof}

\begin{remark}[Conormal superposition]\label{wave:rem:superposition}
In polar coordinates $\eta=r\omega$, homogeneity gives
$\varphi_{\pm}(t,x,r\omega)-ry\cdot\omega=rh_\pm(t,x,y,\omega)$, so each term of the parametrix is
\begin{equation*}
(2\pi)^{1-n}\int_{\Sphere^{n-1}}u_\pm(t,x;y,\omega)\,d\omega,\qquad
u_\pm=\frac1{2\pi}\int_0^\infty e^{irh_\pm(t,x,y,\omega)}a_\pm(t,x,y,r\omega)\chi(r\omega)r^{n-1}\,dr .
\end{equation*}
For fixed $(y,\omega)$, $u_\pm$ is a conormal distribution, of order $\frac{3n-3}4$, associated
with the hypersurface $H_\pm(y,\omega)=\{h_\pm=0\}\subset(-t_0,t_0)\times W_k$, which is
characteristic for $P$ ($(\pa_th_\pm)^2=|d_xh_\pm|_g^2$ by \eqref{wave:eq:eikonal}) and
meets $\{t=0\}$ in the hyperplane $\{(x-y)\cdot\omega=0\}$; there are exactly two such
characteristic hypersurfaces, corresponding to the two points $\tau=\pm|\omega|_g$ of
$\Sigma(P)$ over the conormal $\omega\,dx$. The amplitude of $u_\pm$ solves the transport
equation, and at $t=0$ the sum $u_++u_-$ reproduces the plane wave decomposition
\eqref{wave:eq:plane-waves} of $\delta_y$. As a function of all variables
$(t,x,y,\omega)$, $u_\pm$ is conormal with respect to the hypersurface
$\{h_\pm=0\}$, smooth since $d_yh_\pm=-\omega\ne0$, and the kernel is its push-forward
along $\Sphere^{n-1}$.
\end{remark}

\begin{corollary}\label{wave:cor:WF-short}
For $|t|<t_0$ the wavefront sets of $K_{\cos}$ and $K_{\sin}$ are contained in
\begin{equation*}
\Gamma=\big\{(t,x,y;\tau,\xi,-\eta);\ \eta\ne0,\ \tau=\pm|\eta|_{g(y)},\ (x,\xi)=G_{\mp t}(y,\eta)\big\}.
\end{equation*}
In particular their singular supports are contained in $\{d(x,y)=|t|\}$ for
$|t|<\min(t_0,\inj X)$.
\end{corollary}
\begin{proof}
By \eqref{wave:eq:WF-osc}, the wavefront set of each term is contained in the set
of points
\begin{equation*}
(t,x,y;\pa_t\phi_\pm,d_x\phi_\pm,d_y\phi_\pm)\quad\text{with}\quad d_\eta\phi_\pm=0,\qquad \phi_\pm=\varphi_{k,\pm}(t,x,\eta)-y\cdot\eta .
\end{equation*}
By \eqref{wave:eq:critical} there $(x,d_x\varphi_\pm)=G_{\mp t}(y,\eta)$, while
$d_y\phi_\pm=-\eta$ and $\pa_t\phi_\pm=\pm q(x,d_x\varphi_\pm)=\pm q(y,\eta)$ since $G$
preserves $q$. Within the injectivity radius the geodesic of length $|t|$ is
minimising, so $d(x,y)=|t|$.
\end{proof}

By \cref{wave:lem:bichar}, $\Gamma$ is the union of the null bicharacteristics of
$D_t^2-\Delta_x$ (acting in $(t,x)$, with $(y,\eta)$ fixed) through the points
$(0,y,y;\pm|\eta|,\eta,-\eta)$, i.e.\ the flow-out of the wavefront set of the
initial datum $\delta(x-y)$. In words: the singularity of $\delta_y$ travels out
along all unit speed geodesics from $y$, and at time $t$ it sits on the geodesic
sphere of radius $|t|$ about $y$ (\cref{wave:fig:wavefront}).

\begin{figure}[ht]
  \centering
  \begin{tikzpicture}[>=Stealth,scale=0.95]
    \fill[gray!12] (-2.5,-0.7) rectangle (2.5,0.7);
    \draw[->] (-2.8,0) -- (2.9,0) node[right]{$x$};
    \draw[->] (0,-0.9) -- (0,5.4) node[above]{$t$};
    \draw (-2.5,-0.9) -- (-2.5,5.2) (2.5,-0.9) -- (2.5,5.2);
    \node[below,font=\small] at (-2.5,-0.9) {$y-\frac L2$};
    \node[below,font=\small] at (2.5,-0.9) {$y+\frac L2$};
    \node[below,font=\small] at (0,-0.9) {$y$};
    \draw[very thick] (0,0) -- (2.5,2.5) (-2.5,2.5) -- (0,5);
    \draw[very thick] (0,0) -- (-2.5,2.5) (2.5,2.5) -- (0,5);
    \draw[very thick] (0,0) -- (-0.7,-0.7) (0,0) -- (0.7,-0.7);
    \draw[dotted] (-2.5,5) -- (2.5,5);
    \node[right,font=\small] at (2.6,5) {$t=L$: $x=y$ again};
    \node[right,font=\small] at (2.6,2.5) {$t=\frac L2$};
    \node[right,font=\small] at (2.6,0.35) {$|t|<t_0$};
    \node[above left,font=\small] at (2.0,2.0) {$x=y+t$};
  \end{tikzpicture}
  \caption{The singular support of $K_{\cos}(t,\cdot,y)$ on the circle
  $X=\R/L\Z$, drawn on a fundamental domain $|x-y|\le\frac L2$ whose ends are
  identified. For small $|t|$ (shaded) it is the geodesic sphere $\{d(x,y)=|t|\}$, as
  in \cref{wave:cor:WF-short}. For larger $t$ the singularity keeps moving along the
  geodesics and returns to $x=y$ at $t=L$, the length of the closed geodesic; this is
  why the wave trace is singular at $t=L$ (\cref{wave:sec:poisson}).}
  \label{wave:fig:wavefront}
\end{figure}

\section{Operations on conormal distributions}\label{wave:sec:operations}

We record the behaviour of conormal distributions under push-forward along
fibres and restriction to transversal submanifolds, and the reduction of
amplitudes depending on the normal variable. Together these compute the trace
of the wave group in \cref{wave:sec:trace}.

\begin{proposition}[Push-forward]\label{wave:prop:pushforward}
Let $M_1$ be a manifold of dimension $N_1$, $Z$ a compact manifold of dimension $d$
with a smooth positive density $\nu$, $\pi\colon M_1\times Z\to M_1$ the projection, and
$H_1\subset M_1$ a hypersurface. If $u\in\Ic{m}(M_1\times Z,H_1\times Z)$ then
$\pi_*u=\int_Zu\,\nu\in\Ic{m+\frac d4}(M_1,H_1)$; in local coordinates, with
$u=\frac1{2\pi}\int e^{is\rho}a(z',w,\rho)\,d\rho$, the push-forward has amplitude
$\int_Za(z',w,\rho)\,\nu(w)$.
\end{proposition}
\begin{proof}
$\pi_*$ maps $\CI$ to $\CI$ and distributions smooth off $H_1\times Z$ to
distributions smooth off $H_1$. Near $H_1$ choose coordinates $(s,z')$ on $M_1$ with
$H_1=\{s=0\}$; then $s\circ\pi$ defines $H_1\times Z$ and, using a partition of unity
on $Z$ and coordinates $w$ there, $u$ is locally as in \eqref{wave:eq:conormal-local}
with $a(z',w,\rho)\in\Sym{\mu}$, $\mu=m+\frac{N_1+d}4-\frac12$. Integrating in $w$
gives an amplitude in $\Sym{\mu}(\R^{N_1-1};\R)$, and $\mu=m'+\frac{N_1}4-\frac12$ with
$m'=m+\frac d4$.
\end{proof}

\begin{proposition}[Restriction]\label{wave:prop:restriction}
Let $H\subset M$ be a hypersurface, $\dim M=N$, and $Y\subset M$ a submanifold of
codimension $k$ transversal to $H$ ($T_zH+T_zY=T_zM$ for $z\in H\cap Y$). Then
restriction extends from $\CI$ to a map
$\Ic{m}(M,H)\to\Ic{m+\frac k4}(Y,H\cap Y)$, and the local amplitude restricts.
\end{proposition}
\begin{proof}
By transversality there are local coordinates $(s,z',z'')$ with $H=\{s=0\}$ and
$Y=\{z''=0\}$, $z''\in\R^k$. Then $u=\frac1{2\pi}\int e^{is\rho}a(z',z'',\rho)\,d\rho$ depends
smoothly on $z''$ as a distribution in $(s,z')$, and $u|_Y=\frac1{2\pi}\int e^{is\rho}a(z',0,\rho)\,d\rho$,
consistent with the restriction of smooth functions (and with $N^*Y\cap N^*H=0$).
The symbol order $\mu=m+\frac N4-\frac12$ is unchanged, and
$\mu=m'+\frac{N-k}4-\frac12$ gives $m'=m+\frac k4$.
\end{proof}

\begin{lemma}\label{wave:lem:fourier-amplitude}
Let $B\in\CI(\R\times\R)$ satisfy $B(t,\rho)=0$ for $|t|\ge t_1$ and
$|\pa_t^k\pa_\rho^lB(t,\rho)|\le C_{kl}\ang{\rho}^{m-l}$. Then
\begin{equation*}
J_B(\lambda)=\int_\R\Big(\int_\R e^{it(\rho-\lambda)}B(t,\rho)\,dt\Big)d\rho
\end{equation*}
(an absolutely convergent iterated integral) lies in $\Sym{m}(\R)$, and
$J_B(\lambda)\sim2\pi\sum_{k\ge0}\frac{i^k}{k!}\pa_t^k\pa_\rho^kB(0,\lambda)$; in particular
$J_B(\lambda)-2\pi B(0,\lambda)\in\Sym{m-1}(\R)$. If $B_\eps$ is bounded in this class
for $\eps\in(0,1]$ and $B_\eps\to B$ pointwise, then $J_{B_\eps}\to J_B$ pointwise.
Consequently $v=\frac1{2\pi}\int e^{it\rho}B(t,\rho)\,d\rho\in\Ic{m+\frac14}(\R,\{0\})$, with
$\hat v=J_B/(2\pi)$.
\end{lemma}
\begin{proof}
Let $\hat B(\sigma,\rho)=\int e^{-it\sigma}B(t,\rho)\,dt$; then
$|\pa_\sigma^j\pa_\rho^l\hat B|\le C_{Njl}\ang{\sigma}^{-N}\ang{\rho}^{m-l}$, and
$J_B(\lambda)=\int\hat B(\lambda-\rho,\rho)\,d\rho=\int\hat B(\sigma,\lambda-\sigma)\,d\sigma$. Taylor's
formula in the second variable gives
$\hat B(\sigma,\lambda-\sigma)=\sum_{k<K}\frac{(-\sigma)^k}{k!}\pa_\rho^k\hat B(\sigma,\lambda)+R_K$ with,
by Peetre's inequality $\ang{\lambda-s\sigma}^{\mu}\le2^{|\mu|}\ang{\lambda}^{\mu}\ang{\sigma}^{|\mu|}$,
$|R_K|\le C_N\ang{\sigma}^{K+|m-K|-N}\ang{\lambda}^{m-K}$, which integrates to
$O(\ang{\lambda}^{m-K})$. Since $\sigma^k\hat B=\widehat{D_t^kB}$ and $\int\hat f\,d\sigma=2\pi f(0)$, the
$k$-th term is $2\pi\frac{(-1)^k}{k!}D_t^k\pa_\rho^kB(0,\lambda)=2\pi\frac{i^k}{k!}\pa_t^k\pa_\rho^kB(0,\lambda)$.
Finally $\pa_\lambda J_B=J_{\pa_\rho B}$ and $\pa_\rho B$ is of the same type with $m$
replaced by $m-1$, so the bound $|J_B|\le C\ang{\lambda}^m$ (case $K=0$) gives the symbol
estimates. The convergence statement follows by dominated convergence from the
uniform bound $|\hat B_\eps(\lambda-\rho,\rho)|\le C\ang{\lambda-\rho}^{-N}\ang{\rho}^m$. For the last
statement, $\hat v(\lambda)=\frac1{2\pi}J_B(\lambda)\in\Sym{m}$ and $v$ has compact support,
so $v\in\Ic{m+\frac14}(\R,\{0\})$.
\end{proof}

In \cref{wave:lem:fourier-amplitude} the amplitude $B(t,\rho)$ is replaced by the
amplitude $\hat v(\rho)$ independent of $t$; this is the one-dimensional case of
the statement that amplitudes may depend on the normal variable.

\section{The trace of the wave group}\label{wave:sec:trace}

We now integrate the parametrix over the diagonal. Only the restriction of the
kernel to $\{x=y\}$ matters, and for small $|t|\ne0$ this restriction is smooth,
because the singularities sit on $\{d(x,y)=|t|\}$; all the singularity of the trace
near $t=0$ therefore comes from $t=0$ itself.

For $\psi\in\cS(\R)$ let
$C_\psi=\int\psi(t)\cos(t\Lambda)\,dt=\tilde\psi(\Lambda)$, where
$\tilde\psi(\lambda)=\int\psi(t)\cos(t\lambda)\,dt=\frac12(\hat\psi(\lambda)+\hat\psi(-\lambda))$ is rapidly
decreasing. By \cref{wave:lem:apriori} and the polynomial bounds on
$\norm{\phi_j}_{\mathcal C^k}$, the kernel
$K_{C_\psi}(x,y)=\sum_j\tilde\psi(\lambda_j)\phi_j(x)\phi_j(y)$ converges in $\CI(X\times X)$;
$C_\psi$ is smoothing and of trace class, and
\begin{equation*}
\Tr C_\psi=\sum_j\tilde\psi(\lambda_j)=\int_XK_{C_\psi}(x,x)\dvol(x).
\end{equation*}
Since $\sum_j\ang{\lambda_j}^{-M}<\infty$ for $M>n$, this is a continuous linear functional
on $\cS(\R)$:

\begin{definition}\label{wave:def:trace}
The \emph{(even) wave trace} is the tempered distribution
\begin{equation*}
T(t)=\Tr\cos(t\Lambda)=\sum_{j\ge0}\cos(t\lambda_j),\qquad T(\psi)=\Tr C_\psi .
\end{equation*}
\end{definition}

The trace of the wave group $U(t)$ is $2T$. With the measure $\mu=\sum_j\delta_{\lambda_j}$
on $[0,\infty)$ and its even part $\mu_e=\frac12\sum_j(\delta_{\lambda_j}+\delta_{-\lambda_j})$,
$T(t)=\int e^{it\sigma}d\mu_e(\sigma)$, i.e.\ $\hat T=2\pi\mu_e$, and $N(\lambda)=\mu([0,\lambda])$. For
$\psi\in\CIc(\R)$,
\begin{equation}\label{wave:eq:psiT-hat}
\widehat{\psi T}(\lambda)=T(\psi e^{-i\lambda\cdot})=\frac12\sum_j\big(\hat\psi(\lambda-\lambda_j)+\hat\psi(\lambda+\lambda_j)\big).
\end{equation}
Let $\omega_n$ denote the volume of the unit ball in $\R^n$ (so $\vol(\Sphere^{n-1})=n\omega_n$).

\begin{theorem}\label{wave:thm:trace}
Let $t_0>0$ be as in \cref{wave:thm:parametrix}. For every $\psi\in\CIc((-t_0,t_0))$,
$\psi T\in\Ic{n-\frac34}(\R,\{0\})$, i.e.\ $\widehat{\psi T}\in\Sym{n-1}(\R)$; more precisely,
$\widehat{\psi T}$ is polyhomogeneous as $\lambda\to\pm\infty$ and
\begin{equation}\label{wave:eq:trace-leading}
\widehat{\psi T}(\lambda)=\frac{n\,\omega_n}{2(2\pi)^{n-1}}\,\psi(0)\vol(X)\,|\lambda|^{n-1}+O(|\lambda|^{n-2}),\qquad|\lambda|\to\infty .
\end{equation}
\end{theorem}
\begin{proof}
\emph{Reduction to the parametrix.} With $\psi_\lambda=\psi e^{-i\lambda\cdot}$,
$\widehat{\psi T}(\lambda)=\int_XK_{C_{\psi_\lambda}}(x,x)\dvol$ (the definition of $C_\psi$
extends to complex $\psi$), and $K_{C_{\psi_\lambda}}=\int\psi_\lambda(t)K_{\cos}(t)\,dt$. By
\cref{wave:thm:parametrix}, $K_{\cos}=\sum_{k,\pm}S_{k,\pm}+R$, with $S_{k,\pm}$ the terms of the parametrix, on
$(-t_0,t_0)\times X\times X$ with $R$ smooth; the contribution
$\iint\psi(t)e^{-it\lambda}R(t,x,x)\,dt\dvol$ of $R$ is $O(|\lambda|^{-\infty})$. Fix $k$ and
the sign $+$ (the sign $-$ is analogous) and write $\varphi=\varphi_{k,+}$,
$a=a_{k,+}$. Then
\begin{gather*}
\int\psi_\lambda(t)S_{k,+}(t,x,y)\,dt=(2\pi)^{-n}\int e^{-iy\cdot\eta}\alpha(x,y,\eta)\,d\eta,\\
\alpha=\zeta_k(x)\chi_k(y)\chi(\eta)\int e^{i(\varphi(t,x,\eta)-t\lambda)}\psi(t)a\,dt .
\end{gather*}
For fixed $\lambda$, $|\pa_t\varphi-\lambda|=|q(x,d_x\varphi)-\lambda|\ge c|\eta|$ for $|\eta|\ge C(1+|\lambda|)$, so
repeated integration by parts with $(\pa_t\varphi-\lambda)^{-1}D_t$ shows that $\alpha$ and all its
$(x,y)$-derivatives are rapidly decreasing in $\eta$. Hence this kernel is smooth, given
by an absolutely convergent integral, and its value at $x=y$ is the absolutely
convergent iterated integral $(2\pi)^{-n}\int(\int\cdots dt)\,d\eta$. Thus
\begin{equation*}
I_{k,+}(\lambda)=(2\pi)^{-n}\int_X\int_{\R^n}\Big(\int_\R e^{i(\varphi(t,x,\eta)-x\cdot\eta-t\lambda)}b(t,x,\eta)\,dt\Big)d\eta\dvol(x),
\end{equation*}
with $b(t,x,\eta)=\psi(t)\chi_k(x)a(t,x,x,\eta)\chi(\eta)$ (using $\zeta_k\chi_k=\chi_k$). By
dominated convergence we may insert $\chi_0(\eps\eta)$, $\chi_0\in\CIc(\R^n)$, $\chi_0(0)=1$,
and let $\eps\to0$; for $\eps>0$ the integral is absolutely convergent and Fubini's
theorem allows the $t$-integration to be done last.

\emph{Restriction to the diagonal.} Since $\varphi(0,x,\eta)=x\cdot\eta$,
\begin{equation*}
\varphi(t,x,\eta)-x\cdot\eta=tQ(t,x,\eta),\qquad Q(t,x,\eta)=\int_0^1q\big(x,d_x\varphi(st,x,\eta)\big)\,ds,
\end{equation*}
with $Q$ smooth, positive, homogeneous of degree one in $\eta$ and
$Q(0,x,\eta)=|\eta|_{g(x)}$. (Thus the restriction of the characteristic hypersurface
$\{h_+=0\}$ of \cref{wave:rem:superposition} to the diagonal is exactly $\{t=0\}$, the
transversal situation of \cref{wave:prop:restriction}.) In polar coordinates
$\eta=r\omega$ and then $\rho=rQ(t,x,\omega)$ for fixed $(t,x,\omega)$,
$d\eta=\rho^{n-1}Q(t,x,\omega)^{-n}d\rho\,d\omega$, and we obtain
$I_{k,+}(\lambda)=\lim_{\eps\to0}\iint e^{it(\rho-\lambda)}B_\eps(t,\rho)\,dt\,d\rho$ with
\begin{equation*}
B_\eps(t,\rho)=(2\pi)^{-n}\int_X\int_{\Sphere^{n-1}}b\Big(t,x,\frac{\rho\omega}{Q}\Big)\chi_0\Big(\frac{\eps\rho\omega}{Q}\Big)\frac{\rho^{n-1}}{Q(t,x,\omega)^n}\,d\omega\dvol(x)
\end{equation*}
for $\rho>0$, and $B_\eps=0$ for $\rho\le0$ (as $\chi=0$ near $0$). This is the
push-forward of \cref{wave:prop:pushforward} along $X\times\Sphere^{n-1}$. The
$B_\eps$ are bounded in the class of \cref{wave:lem:fourier-amplitude} with $m=n-1$ and
converge to $B=B_0$, whose amplitude is polyhomogeneous in $\rho$. By
\cref{wave:lem:fourier-amplitude}, $I_{k,+}=J_B\in\Sym{n-1}$, polyhomogeneous, and
$I_{k,+}(\lambda)=2\pi B(0,\lambda)+O(\ang{\lambda}^{n-2})$. In particular $I_{k,+}(\lambda)=O(|\lambda|^{-\infty})$
as $\lambda\to-\infty$.

\emph{Leading term.} At $t=0$, $b(0,x,\eta)=\psi(0)\chi_k(x)a(0,x,x,\eta)$ for $|\eta|\ge1$, and by
\eqref{wave:eq:initial-symbols} $a(0,x,x,\eta)=\frac12|g(x)|^{-1/2}+\Sym{-1}$. Hence for
$\rho\to+\infty$
\begin{equation*}
B(0,\rho)=\frac{\psi(0)}{2(2\pi)^n}\rho^{n-1}\int_X\chi_k(x)|g(x)|^{-1/2}\int_{\Sphere^{n-1}}|\omega|_{g(x)}^{-n}\,d\omega\dvol(x)+O(\rho^{n-2}).
\end{equation*}
In polar coordinates the ellipsoid $\{\eta;\ |\eta|_{g(x)}\le1\}$ has volume
$\frac1n\int_{\Sphere^{n-1}}|\omega|_{g(x)}^{-n}d\omega$, and it is the image of the unit ball under a linear
map of determinant $\det(g^{ij})^{-1/2}=|g(x)|^{1/2}$; so the inner integral is $n\omega_n|g(x)|^{1/2}$ and
$B(0,\rho)=\frac{n\omega_n}{2(2\pi)^n}\psi(0)\rho^{n-1}\int_X\chi_k\dvol+O(\rho^{n-2})$. The terms with the sign
$-$ have phase $\varphi_{k,-}-x\cdot\eta=-tQ_-$ and give the same expression with $\lambda$ replaced by
$-\lambda$, rapidly decreasing as $\lambda\to+\infty$. Summing over $k$ ($\sum\chi_k=1$) and
multiplying by $2\pi$ gives \eqref{wave:eq:trace-leading}.
\end{proof}

\begin{remark}\label{wave:rem:circle}
For the circle $X=\R/L\Z$, $\lambda=2\pi|k|/L$ for $k\in\Z$ and
$T(t)=\sum_{k\in\Z}\cos(2\pi kt/L)$, while by Poisson summation $\sum_k\hat\psi(\lambda-2\pi k/L)
=L\sum_m\psi(mL)e^{-imL\lambda}=L\psi(0)$ if $\supp\psi\subset(-L,L)$. Thus
$\widehat{\psi T}=L\psi(0)$, in agreement with \eqref{wave:eq:trace-leading} ($n=1$,
$\omega_1=2$). The restriction on the support of $\psi$ is necessary: $T$ is singular at every
$t\in L\Z$, the lengths of the closed geodesics (\cref{wave:sec:poisson}).
\end{remark}

\section{Weyl asymptotics}\label{wave:sec:weyl}

\begin{discussion}
Weyl's law has a simple heuristic. By the uncertainty principle each eigenfunction
occupies a cell of volume $(2\pi)^n$ in phase space, and an eigenfunction with
$\lambda_j\le\lambda$ ``lives'' in the region $\{(x,\xi);\ |\xi|_{g(x)}\le\lambda\}$ of
$T^*X$. That region has volume $\omega_n\vol(X)\lambda^n$, so one expects
$N(\lambda)\approx(2\pi)^{-n}\omega_n\vol(X)\lambda^n$. The proof below makes this
precise with the trace formula of \cref{wave:thm:trace}, whose leading coefficient is
exactly this phase-space volume, written as a density in $\lambda$.
\end{discussion}

\begin{theorem}[Weyl's law with remainder]\label{wave:thm:weyl}
Let $X$ be a compact Riemannian manifold of dimension $n\ge2$. Then
\begin{equation*}
N(\lambda)=(2\pi)^{-n}\omega_n\vol(X)\lambda^n+O(\lambda^{n-1}),\qquad\lambda\to\infty .
\end{equation*}
\end{theorem}
\Cref{wave:fig:weyl} compares $N(\lambda)$ with $C_W\lambda^n$ for the flat torus
$\R^2/2\pi\Z^2$, where $N(\lambda)$ counts lattice points in a disc
(\cref{wave:ex:torus}).

\begin{figure}[ht]
  \centering
  \begin{tikzpicture}[>=Stealth]
    \draw[->] (0,0) -- (9.5,0) node[right]{$\lambda$};
    \draw[->] (0,0) -- (0,6.0) node[above]{$N(\lambda)$};
    \foreach \l in {1,...,6}{\draw (1.5*\l,0.08) -- (1.5*\l,-0.08) node[below,font=\footnotesize]{$\l$};}
    \foreach \n in {25,50,75,100}{\draw (0.08,0.05*\n) -- (-0.08,0.05*\n) node[left,font=\footnotesize]{$\n$};}
    \draw[thick,dashed,domain=0:6,samples=60,smooth,variable=\l] plot ({1.5*\l},{0.05*pi*\l*\l});
    \draw[thick] (0.000,0.000) -- (0.000,0.050) -- (1.500,0.050) -- (1.500,0.250) -- (2.121,0.250) -- (2.121,0.450) -- (3.000,0.450) -- (3.000,0.650) -- (3.354,0.650) -- (3.354,1.050) -- (4.243,1.050) -- (4.243,1.250) -- (4.500,1.250) -- (4.500,1.450) -- (4.743,1.450) -- (4.743,1.850) -- (5.408,1.850) -- (5.408,2.250) -- (6.000,2.250) -- (6.000,2.450) -- (6.185,2.450) -- (6.185,2.850) -- (6.364,2.850) -- (6.364,3.050) -- (6.708,3.050) -- (6.708,3.450) -- (7.500,3.450) -- (7.500,4.050) -- (7.649,4.050) -- (7.649,4.450) -- (8.078,4.450) -- (8.078,4.850) -- (8.485,4.850) -- (8.485,5.050) -- (8.746,5.050) -- (8.746,5.450) -- (9.000,5.450) -- (9.000,5.650);
    \node[font=\small] at (5.2,4.3) {$C_W\lambda^2=\pi\lambda^2$};
    \draw[->,thin] (5.6,4.05) -- (6.55,3.45);
  \end{tikzpicture}
  \caption{The counting function $N(\lambda)=\#\{k\in\Z^2;\ |k|\le\lambda\}$ of the flat
  torus $\R^2/2\pi\Z^2$ (staircase) and the Weyl term $C_W\lambda^2=\pi\lambda^2$
  (dashed). The jumps of $N$ are the multiplicities, here at most of order
  $\lambda^{\epsilon}$; on the round sphere they are of order $\lambda^{n-1}$, which
  is why the remainder $O(\lambda^{n-1})$ cannot be improved in general.}
  \label{wave:fig:weyl}
\end{figure}

The proof has three steps. The trace formula controls a \emph{smoothed} spectral
density $\rho*\mu$, where $\rho$ is a fixed positive bump. Positivity of $\rho$ gives
an a priori bound on the number of eigenvalues in any window of fixed length
(Step 1). Integrating $\rho*\mu$ gives a smoothed counting function with the correct
asymptotics (Step 2), and the a priori bound shows that smoothing changes $N$ by at
most $O(\lambda^{n-1})$ (Step 3). \Cref{wave:fig:tauberian} shows the smoothing.

\begin{figure}[ht]
  \centering
  \begin{tikzpicture}[>=Stealth]
    \draw[->] (0,0) -- (9.5,0) node[right]{$\lambda$};
    \draw[->] (0,0) -- (0,5.3);
    \foreach \l in {1,...,6}{\draw (1.5*\l,0.08) -- (1.5*\l,-0.08) node[below,font=\footnotesize]{$\l$};}
    \foreach \x/\h in {0.000/0.200, 1.500/0.800, 2.121/0.800, 3.000/0.800, 3.354/1.600, 4.243/0.800, 4.500/0.800, 4.743/1.600, 5.408/1.600, 6.000/0.800, 6.185/1.600, 6.364/0.800, 6.708/1.600, 7.500/2.400, 7.649/1.600, 8.078/1.600, 8.485/0.800, 8.746/1.600, 9.000/0.800}
      {\draw[gray,line width=1.2pt] (\x,0) -- (\x,\h);}
    \draw[thick,dashed] (0,0) -- (9,4.524) node[right,font=\small]{$2\pi\lambda$};
    \draw[very thick] plot[smooth] coordinates {(0.000,0.146) (0.150,0.152) (0.300,0.158) (0.450,0.172) (0.600,0.205) (0.750,0.265) (0.900,0.353) (1.050,0.466) (1.200,0.594) (1.350,0.721) (1.500,0.832) (1.650,0.918) (1.800,0.973) (1.950,1.003) (2.100,1.021) (2.250,1.046) (2.400,1.097) (2.550,1.187) (2.700,1.316) (2.850,1.466) (3.000,1.610) (3.150,1.718) (3.300,1.773) (3.450,1.781) (3.600,1.765) (3.750,1.765) (3.900,1.811) (4.050,1.916) (4.200,2.070) (4.350,2.239) (4.500,2.387) (4.650,2.487) (4.800,2.531) (4.950,2.533) (5.100,2.519) (5.250,2.521) (5.400,2.563) (5.550,2.659) (5.700,2.802) (5.850,2.972) (6.000,3.136) (6.150,3.261) (6.300,3.323) (6.450,3.323) (6.600,3.285) (6.750,3.253) (6.900,3.273) (7.050,3.371) (7.200,3.540) (7.350,3.744) (7.500,3.929) (7.650,4.050) (7.800,4.089) (7.950,4.062) (8.100,4.012) (8.250,3.988) (8.400,4.026) (8.550,4.136) (8.700,4.304) (8.850,4.495) (9.000,4.670)};
    \node[font=\small] at (2.6,2.6) {$\rho*\mu$};
    \draw[->,thin] (2.6,2.35) -- (3.1,1.8);
  \end{tikzpicture}
  \caption{The Tauberian step for the torus of \cref{wave:fig:weyl}. The spectral
  measure $\mu=\sum_j\delta_{\lambda_j}$ (gray spikes, heights proportional to
  multiplicities) is too irregular to have asymptotics, but its convolution with a
  fixed positive bump $\rho$ of width about $\frac13$ (thick curve) follows the
  smooth density $nC_W\lambda^{n-1}=2\pi\lambda$ (dashed) up to an error of lower
  order.}
  \label{wave:fig:tauberian}
\end{figure}

\begin{proof}
Write $C_W=(2\pi)^{-n}\omega_n\vol(X)$.

\emph{Choice of a test function.} Let $\theta\in\CIc((-t_0/2,t_0/2))$ be real, even,
non-negative and not identically zero, normalised by $\int\theta^2=1$, and let
$\psi=\theta*\theta$. Then $\psi\in\CIc((-t_0,t_0))$ is even, $\psi(0)=1$ and
$\hat\psi=\hat\theta^2=|\hat\theta|^2\ge0$. Set $\rho=\hat\psi/(2\pi)$: $\rho\in\cS(\R)$ is even and
non-negative, $\int\rho=\psi(0)=1$, and $\rho(0)=(2\pi)^{-1}(\int\theta)^2>0$, so there are
$a,c_0>0$ with $\rho\ge c_0$ on $[-a,a]$.

\emph{Smoothed density.} By \eqref{wave:eq:psiT-hat},
$(\rho*\mu)(\lambda)=\sum_j\rho(\lambda-\lambda_j)=\pi^{-1}\widehat{\psi T}(\lambda)-\sum_j\rho(\lambda+\lambda_j)$. For
$\lambda\ge0$, $\sum_j\rho(\lambda+\lambda_j)\le C_M\sum_j(1+\lambda+\lambda_j)^{-M}=O(\lambda^{-\infty})$, by
\cref{wave:lem:apriori}. Hence by \cref{wave:thm:trace}
\begin{equation}\label{wave:eq:smoothed-density}
(\rho*\mu)(\lambda)=nC_W\lambda^{n-1}+O(\lambda^{n-2}),\quad\lambda\to+\infty;\qquad(\rho*\mu)(\lambda)=O(\ang{\lambda}^{-\infty}),\quad\lambda\to-\infty,
\end{equation}
the second because $\lambda-\lambda_j\le\lambda<0$. Note $\pi^{-1}\frac{n\omega_n}{2(2\pi)^{n-1}}=n(2\pi)^{-n}\omega_n$.

\emph{Step 1: a priori bound.} Since $\rho\ge0$ and $\rho\ge c_0$ on $[-a,a]$,
$\mu([\lambda-a,\lambda+a])\le c_0^{-1}(\rho*\mu)(\lambda)\le C(1+|\lambda|)^{n-1}$ for all $\lambda\in\R$
(the function $\rho*\mu$ is continuous). Covering an interval of length $2r$ by at most
$1+r/a$ intervals of length $2a$,
\begin{equation}\label{wave:eq:local-bound}
\mu([\lambda-r,\lambda+r])\le C'(1+r)^n(1+|\lambda|)^{n-1},\qquad\lambda\in\R,\ r\ge0 .
\end{equation}

\emph{Step 2: the smoothed counting function.} Let
$N_\rho(\lambda)=\int N(\lambda-\sigma)\rho(\sigma)\,d\sigma=\sum_jR(\lambda-\lambda_j)$, $R(s)=\int_{-\infty}^s\rho$.
Termwise differentiation (justified by the locally uniform convergence of
$\sum_j\rho(\lambda-\lambda_j)$) gives $N_\rho'=\rho*\mu$, and $N_\rho(\lambda)\to0$ as $\lambda\to-\infty$.
Thus, by \eqref{wave:eq:smoothed-density} and $n\ge2$,
\begin{equation*}
N_\rho(\lambda)=\int_{-\infty}^1(\rho*\mu)\,ds+\int_1^\lambda\big(nC_Ws^{n-1}+O(s^{n-2})\big)\,ds=C_W\lambda^n+O(\lambda^{n-1}).
\end{equation*}

\emph{Step 3: removing the smoothing.} Since $\int\rho=1$, $\rho\ge0$ and $N$ is
non-decreasing, $|N(\lambda)-N(\lambda-\sigma)|\le\mu([\lambda-|\sigma|,\lambda+|\sigma|])$, so by
\eqref{wave:eq:local-bound}
\begin{equation*}
|N(\lambda)-N_\rho(\lambda)|\le\int|N(\lambda)-N(\lambda-\sigma)|\rho(\sigma)\,d\sigma
\le C'(1+|\lambda|)^{n-1}\int(1+|\sigma|)^n\rho(\sigma)\,d\sigma .
\end{equation*}
Combining with Step 2 proves the theorem.
\end{proof}

\begin{remark}\label{wave:rem:weyl-sharp}
\begin{enumerate}[label=(\roman*)]
\item The remainder cannot be improved in general: on the round sphere
  $\Sphere^n$ the eigenvalues $k(k+n-1)$ of $\Delta$ have multiplicities of order $k^{n-1}$,
  so $N$ has jumps of order $\lambda^{n-1}$ (\cref{wave:ex:sphere}). The remainder is
  $o(\lambda^{n-1})$ if the set of periodic points of the geodesic flow on the unit
  cosphere bundle has measure zero (Duistermaat--Guillemin, Ivrii); this uses the
  structure of $T$ at the non-zero periods. The weaker formula
  $N(\lambda)\sim C_W\lambda^n$ can also be obtained from the heat kernel or, as in
  \cite{Guillemin-Weyl}, from residues of the zeta function of $\Delta$.
\item For $n=1$ each component of $X$ is a circle and everything is explicit
  (\cref{wave:rem:circle}); in that case $N(\lambda)=C_W\lambda+O(1)$, but the argument of Step 2
  only gives $O(\log\lambda)$ unless one uses that the lower order terms in
  \eqref{wave:eq:trace-leading} vanish.
\item The proof only used: the trace formula \eqref{wave:eq:trace-leading} for
  \emph{one} positive test function with small support, the positivity of $\mu$, and a
  polynomial a priori bound. See \cite[Section~17.5]{Hormander-III} for the same
  method based on the Hadamard parametrix.
\end{enumerate}
\end{remark}

\begin{question}
Why must the test function be chosen with $\hat\psi\ge0$?
\end{question}
\begin{answer}
Because Step 1 bounds the number of eigenvalues in a window $[\lambda-a,\lambda+a]$ by
$c_0^{-1}(\rho*\mu)(\lambda)$, which needs $\rho=\hat\psi/(2\pi)\ge0$ everywhere and
$\rho\ge c_0>0$ on $[-a,a]$. With a kernel taking negative values, a large cluster of
eigenvalues near $\lambda$ could be cancelled in $\rho*\mu$ by the contributions of
eigenvalues elsewhere, and the trace formula would give no control on $N$. Since
$\psi$ must also have small support (the parametrix is only valid for $|t|<t_0$),
$\psi=\theta*\theta$ is the natural choice: compact support and $\hat\psi=\hat\theta^2\ge0$.
\end{answer}

\section{Propagation of singularities and the Poisson relation}\label{wave:sec:poisson}

In this section we explain, with some steps only sketched, how the global
behaviour of $T$ is governed by closed geodesics.

\begin{discussion}
The trace $T$ is obtained by restricting the wave kernel to the diagonal $x=y$ and
integrating. A singularity of $K_{\cos}(t,\cdot,y)$ starting at $y$ travels along a
geodesic, so it can be on the diagonal at time $t\ne0$ only if that geodesic returns
to its starting point at time $|t|$ with the same direction, i.e.\ is closed of
length dividing $|t|$ (\cref{wave:fig:wavefront}). After integration in $x$ only these
returns can produce singularities of $T$. Propagation of singularities makes this
argument rigorous for all times.
\end{discussion} The full theory requires Fourier
integral operators \cite{Hormander-FIO} and the propagation theorem for operators of
real principal type, which is proved in \cref{ch:microlocal}; for the strictly hyperbolic case see
\cite[Chapter~XXIII]{Hormander-III}.

Let $C_+\subset\Sigma(P)\times\Sigma(P)$ be the forward bicharacteristic relation:
$(\zeta,\zeta')\in C_+$ if $\zeta$ lies on the null bicharacteristic through $\zeta'$ and
$t(\zeta)\ge t(\zeta')$. Write $\Delta^*$ for the diagonal of
$(T^*(\R\times X)\setminus0)^2$ and $\WFp(E)=\{(\zeta,\zeta');\ (\zeta,-\zeta')\in\WF(E)\}$.

\begin{lemma}\label{wave:lem:WF-E}
$\WFp(E)\cap\{|t-t'|<t_0\}\subset\Delta^*\cup C_+$, and $\WF(E)$ contains no
points with zero covector in either factor.
\end{lemma}
\begin{proof}[Sketch of proof]
By \eqref{wave:eq:E-kernel}, $E=-H(t-t')K_{\sin}(t-t',x,x')$. For
$0<|t-t'|<t_0$, \cref{wave:cor:WF-short} gives $\WF(E)\subset$ pull-back of $\Gamma$, and by
\cref{wave:lem:bichar} a point of $\Gamma$ at time $t-t'>0$ corresponds exactly to a pair
in $C_+$ (the point $(t',x',\tau,\eta)$ is carried to $(t,G_{\mp(t-t')}(x',\eta),\tau)$ for
$\tau=\pm|\eta|$). At $t=t'$, $\WF(H(t-t'))$ consists of covectors $(\sigma,0;-\sigma,0)$ and
$\WF(K_{\sin}(t-t',x,x'))$ of covectors $(\tau,\xi;-\tau,-\xi)$ with $x=x'$, $\tau=\pm|\xi|$; their sum
is never zero, so the product is defined \cite[Section~8.2]{Hormander-I}, and all
resulting points over $\{t=t',x=x'\}$ have the form $(\zeta,-\zeta)$ with $\zeta\ne0$, i.e.\ lie in
$\Delta^*$ in the primed notation. The last claim follows from the same description.
\end{proof}

\begin{theorem}[Propagation of singularities]\label{wave:thm:propagation}
Let $u\in\CmI(\R\times X)$. Then $\WF(u)\setminus\WF(Pu)\subset\Sigma(P)$, and if
$\gamma\colon J\to\Sigma(P)\setminus\WF(Pu)$ is a null bicharacteristic on an interval $J$ with
$\gamma(s_0)\in\WF(u)$ for some $s_0\in J$, then $\gamma(J)\subset\WF(u)$.
\end{theorem}
\begin{proof}
The first statement is microlocal elliptic regularity, since $\Sigma(P)$ is the
characteristic set of $P$. The second follows from the propagation theorem for
operators with real homogeneous principal symbol on Euclidean space proved in
\cref{ch:microlocal}, applied in local coordinates. Let
$A=\{s\in J;\ \gamma(s)\in\WF(u)\}$; it is non-empty, and closed in $J$ since $\WF(u)$
is closed. Let $s_1\in A$. Choose a coordinate chart $\kappa\colon U\to U'\subset\R^{n+1}$ of
$\R\times X$ containing the projection of $\gamma(s_1)$, and a compact interval $I\subset J$
which is a neighbourhood of $s_1$ in $J$ and such that the projection of $\gamma(I)$ lies
in an open set $V$ with $\overline V\subset U$ compact. Let $\phi,\psi\in\CIc(U)$ with
$\phi=1$ near $\overline V$ and $\psi=1$ near $\supp\phi$. In the coordinates,
$P=\sum_{|\alpha|\le2}a_\alpha(y)D_y^\alpha$ with smooth coefficients on $U'$, and
$\tilde P=\sum_{|\alpha|\le2}\psi a_\alpha D_y^\alpha$ (with $\psi$ written in the coordinates)
is a differential operator with coefficients in $\CIc(\R^{n+1})$, hence lies in
$\Psido2(\R^{n+1})$ and has real homogeneous principal symbol $\psi p$. Let
$\tilde u=\phi u$, a distribution with compact support in $U'$. Since $P$ is a
differential operator and $\phi=\psi=1$ on $V$, $\tilde P\tilde u=Pu$ on $V$, so over $V$
the wavefront sets of $\tilde u$ and $\tilde P\tilde u$ coincide with those of $u$ and $Pu$
(the wavefront set is local, and invariantly defined on $T^*(\R\times X)$). Over $V$,
$\psi p=p$, so the Hamilton vector field of $\psi p$ there is the coordinate expression of
$\Ham{p}$, and $\gamma|_I$ is a null bicharacteristic of $\tilde P$ which does not meet
$\WF(\tilde P\tilde u)$ and meets $\WF(\tilde u)$ at $\gamma(s_1)$. By the propagation theorem of
\cref{ch:microlocal}, $\gamma(I)\subset\WF(\tilde u)$, i.e.\ $\gamma(I)\subset\WF(u)$. Hence $A$ is
open in $J$, and since $J$ is connected, $A=J$.
\end{proof}

The same argument applies to the differential operator $D_t^2-\Delta_X$ on $\R\times X\times Y$, whose
principal symbol $\tau^2-|\xi|_g^2$ is real: its null bicharacteristics through points with
$(\tau,\xi)\ne0$ are those of \cref{wave:lem:bichar}, with the $(y,\eta)$ components fixed, while
points with $\tau=0$, $\xi=0$ are stationary (there $\Ham{p}=0$).

\begin{proposition}\label{wave:prop:WF-kernel}
For all $t\in\R$, $\WF(K_{\cos})\subset\Gamma$, with $\Gamma$ as in \cref{wave:cor:WF-short}.
\end{proposition}
\begin{proof}[Sketch of proof]
$K_{\cos}$ is annihilated by $D_t^2-\Delta_x$ and by $D_t^2-\Delta_y$ (as $\cos(t\Lambda)$
commutes with $\Delta$), so $\WF(K_{\cos})\subset\{\tau^2=|\xi|^2=|\eta|^2\}$, which does not meet
$\{\tau=0,\xi=0\}$. By the extension of \cref{wave:thm:propagation}, $\WF(K_{\cos})$ is a
union of complete bicharacteristics of $D_t^2-\Delta_x$; each of these meets
$\{|t|<t_0\}$, where $\WF(K_{\cos})\subset\Gamma$ by \cref{wave:cor:WF-short}. Since $\Gamma$ is a union of
such bicharacteristics, $\WF(K_{\cos})\subset\Gamma$.
\end{proof}

Let $\Lset$ be the set of lengths of closed (non-constant, unit speed) geodesics of
$X$. It is closed, and $\inf\Lset\ge2\inj(X)>0$, since a geodesic loop at $x$ cannot be
shorter than $2\inj(X)$.

\begin{theorem}[Poisson relation]\label{wave:thm:poisson}
$\singsupp T\subset\{0\}\cup\Lset\cup(-\Lset)$.
\end{theorem}
\begin{proof}[Sketch of proof]
Let $\iota\colon\R\times X\to\R\times X\times X$, $(t,x)\mapsto(t,x,x)$, and $\pi\colon\R\times X\to\R$. The
conormal bundle of $\iota(\R\times X)$ consists of covectors $(0,\xi,-\xi)$, which are not in
$\Gamma$ (there $\tau\ne0$), so $\iota^*K_{\cos}$ is defined and
$\WF(\iota^*K_{\cos})\subset\{(t,x;\tau,\xi-\eta);\ (t,x,x;\tau,\xi,-\eta)\in\Gamma\}$
\cite[Section~8.2]{Hormander-I}. Since integration over $t$ against $\psi$ commutes with
$\iota^*$ and turns $K_{\cos}$ into the smooth kernel $K_{C_\psi}$, $T=\pi_*\iota^*K_{\cos}$.
For the push-forward, $\WF(\pi_*v)\subset\{(t,\tau);\ \exists x,\ (t,x;\tau,0)\in\WF(v)\}$. Hence if
$(t,\tau)\in\WF(T)$ there are $x$ and $\eta\ne0$ with $(x,\eta)=G_{\mp t}(x,\eta)$, i.e.\ either
$t=0$, or the geodesic through $x$ in the direction $g^{-1}\eta$ is closed of length dividing
$|t|$ — in particular $|t|\in\Lset$ (a multiple of a closed geodesic is a closed
geodesic).
\end{proof}

\begin{corollary}\label{wave:cor:trace-extended}
The conclusions of \cref{wave:thm:trace} hold for every $\psi\in\CIc((-\inf\Lset,\inf\Lset))$.
\end{corollary}
\begin{proof}
Take $\psi_0\in\CIc((-t_0,t_0))$ with $\psi_0=1$ near $0$. Then $(1-\psi_0)\psi T$ is smooth
with compact support by \cref{wave:thm:poisson}, so its Fourier transform is
rapidly decreasing, and \cref{wave:thm:trace} applies to $\psi_0\psi$, with $(\psi_0\psi)(0)=\psi(0)$.
\end{proof}

\section{Exercises}

\begin{exercise}\label{wave:ex:darboux}
Let $E$ be a real vector space with a non-degenerate antisymmetric bilinear form
$\omega$. Show that $\dim E=2n$ is even and that there is a basis
$v_1,\dots,v_n,w_1,\dots,w_n$ with dual basis $v_j^*,w_j^*$ such that
$\omega=\sum_jv_j^*\wedge w_j^*$. \emph{Hint:} choose $v_1\ne0$, then $w_1$ with
$\omega(v_1,w_1)=1$; then $v_2$ in the $\omega$-orthogonal complement of $\mathrm{span}(v_1,w_1)$ (why
is this complement a complement, and why is $\omega$ non-degenerate on it?), and so on.
Deduce that there is a linear isomorphism $L\colon E\to\R^{2n}$ with $\omega=L^*\omega_0$, where
$\omega_0=\sum_jd\xi_j\wedge dx_j$ is the standard symplectic form.
\end{exercise}

\begin{exercise}\label{wave:ex:lagrangian}
Extend \cref{wave:ex:darboux}: if $V\subset E$ is a Lagrangian subspace ($\dim V=n$,
$\omega|_V=0$), then $L$ can be chosen to map $V$ onto $\{x=0\}\subset\R^{2n}$. \emph{Hint:}
construct the basis with all $v_j\in V$ (or all $w_j$, depending on the sign convention).
\end{exercise}

\begin{exercise}\label{wave:ex:poisson}
Let $(S,\omega)$ be a symplectic manifold and $\Ham{f}$ the Hamilton vector field of
$f\in\CI(S)$, defined by $\iota_{\Ham{f}}\omega=-df$ (or with the opposite sign). Show that the
Poisson bracket $\{f,g\}=\Ham{f}g$ makes $\CI(S)$ into a Lie algebra, and that
$[\Ham{f},\Ham{g}]=\Ham{\{f,g\}}$. \emph{Hint:} the Jacobi identity is equivalent to $d\omega=0$.
\end{exercise}

\begin{exercise}\label{wave:ex:initial-split}
Let $H_0=\{(x-y)\cdot\omega=0\}\subset\R^n$ and let $H_\pm$ be the two characteristic
hypersurfaces of \cref{wave:rem:superposition} through $\{0\}\times H_0$. Show that for
$v_0\in\Ic{m}(\R^n,H_0)$ and $v_1\in\Ic{m+1}(\R^n,H_0)$ there are conormal $u_\pm$,
associated with $H_\pm$ and solving $Pu_\pm\in\CI$ near $t=0$, with $u_++u_-=v_0$ and
$D_t(u_++u_-)=v_1$ at $t=0$ modulo $\CI$, and that their full symbols at $t=0$ are
uniquely determined. (Write the conditions on the symbols as a triangular system,
as in the proof of \cref{wave:thm:parametrix}.)
\end{exercise}

\begin{exercise}\label{wave:ex:volterra}
Let $F\in\Cz([0,t_1]\times X\times X)$ with $|F|\le M$, extended by $0$ for $t<0$, and define the
$t$-convolution $(F\circ G)(t,x,y)=\int_0^t\int_XF(t-s,x,z)G(s,z,y)\dvol(z)\,ds$. Show that the
iterates $F^{(1)}=F$, $F^{(k)}=F^{(k-1)}\circ F$ satisfy
$|F^{(k)}(t,x,y)|\le M^k\vol(X)^{k-1}t^{k-1}/(k-1)!$, and deduce that $\Id+F$ (with $\Id$ the
identity for $t$-convolution, kernel $\delta(t)\delta_y(x)$) has inverse $\Id+\tilde F$,
$\tilde F=\sum_{k\ge1}(-1)^kF^{(k)}$, with $\tilde F$ continuous. Show that if $F$ is smooth and
vanishes to infinite order at $t=0$ then so does $\tilde F$, and use this, together with a local
parametrix, to construct a forward fundamental solution as in \cref{wave:rem:volterra}.
\end{exercise}

\begin{exercise}\label{wave:ex:circle}
For $X=\R/L\Z$ compute $N(\lambda)$ exactly and verify $N(\lambda)=\frac L\pi\lambda+O(1)$. Show that
$T=\sum_{k\in\Z}\cos(2\pi kt/L)=L\sum_{m\in\Z}\delta(t-mL)$ and compare with
\cref{wave:thm:poisson}.
\end{exercise}

\begin{exercise}\label{wave:ex:torus}
For the flat torus $X=\R^n/2\pi\Z^n$ show that the eigenvalues of $\Delta$ are $|k|^2$,
$k\in\Z^n$, and deduce from \cref{wave:thm:weyl} the lattice point estimate
$\#\{k\in\Z^n;\ |k|\le\lambda\}=\omega_n\lambda^n+O(\lambda^{n-1})$ ($n\ge2$).
\end{exercise}

\begin{exercise}\label{wave:ex:sphere}
On the round sphere $\Sphere^2$ the eigenvalues of $\Delta$ are $k(k+1)$, $k\in\N_0$, with
multiplicity $2k+1$. Show that $C_W=1$ and that $N(\lambda)-\lambda^2$ takes values $\ge\frac12\lambda$ and
$\le-\frac12\lambda$ for arbitrarily large $\lambda$, so the remainder in \cref{wave:thm:weyl} is not
$o(\lambda^{n-1})$ in general. Compute $T$ for $\Sphere^2$ and locate its singularities.
\end{exercise}

\begin{exercise}\label{wave:ex:fn}
Show that $f_n$ is homogeneous of degree $-n$, that $f_n=D_s^{n-1}\delta$ for $n$ odd, and that
for $n$ even $f_n(s)=c_ns^{-n}$ for $s\ne0$, with a constant $c_n\ne0$ (compute it). Deduce from
\cref{wave:prop:radon} that, for $n$ odd, $\delta$ is a superposition
$\int_{\Sphere^{n-1}}g_\omega(x\cdot\omega)\,d\omega$ of distributions supported on the
hyperplanes $\{x\cdot\omega=0\}$ (i.e.\ $\supp g_\omega\subset\{0\}$), whereas for $n$ even it is not.
\end{exercise}

\section*{Notes}

Compared with the original notes, we have restricted attention to compact $X$
(dropping the case of $\R^n$ with a metric that is Euclidean at infinity), and solved
the Cauchy problem and constructed the forward fundamental solution by spectral
decomposition, keeping the Volterra-series construction as a remark and an exercise.
The short-time parametrix is written as an oscillatory integral whose phase solves
the eikonal equation; in polar coordinates it is exactly the superposition of
conormal solutions associated with characteristic hypersurfaces through hyperplanes
that the original uses (\cref{wave:rem:superposition}). The order of $f_n(x\cdot\omega)$ in
\cref{wave:lem:fn} is $\frac n2-\frac14$ (not $\frac n4-\frac14$), and Weyl's formula with the
remainder $O(\lambda^{n-1})$ is what the Tauberian argument proves; the original stated
$o(\lambda^{n-1})$, which is false in general (\cref{wave:ex:sphere}).

The remainder $O(\lambda^{n-1})$ is due to Avakumovi\'c, Levitan and H\"ormander
(L.~H\"ormander, \emph{The spectral function of an elliptic operator}, Acta Math.\ 121
(1968)); see also \cite[Section~17.5]{Hormander-III}. The Poisson relation and the
improvement to $o(\lambda^{n-1})$ under a condition on periodic geodesics are due to
Chazarain, Colin de Verdi\`ere and Duistermaat--Guillemin (J.~J.~Duistermaat and
V.~W.~Guillemin, \emph{The spectrum of positive elliptic operators and periodic
bicharacteristics}, Invent.\ Math.\ 29 (1975)), and Ivrii. For Fourier integral operators
and Lagrangian distributions, of which the wave kernel is the basic example, see
\cite{Hormander-FIO}, \cite{Bruning-Guillemin}; for a different proof of the leading term
see \cite{Guillemin-Weyl}.
